% Options for packages loaded elsewhere
\PassOptionsToPackage{unicode}{hyperref}
\PassOptionsToPackage{hyphens}{url}
%
\documentclass[
  11pt,
  a4paper,
]{amsart}
\providecommand{\tightlist}{\setlength{\itemsep}{0pt}\setlength{\parskip}{0pt}}
\usepackage{amsthm}
\usepackage{amsfonts}
\usepackage{mathrsfs}
\usepackage{mathtools}
\usepackage{stmaryrd}
\usepackage{bbm}
\usepackage{enumitem}
\usepackage{slashed}
\usepackage{tikz-cd}
\usepackage{tikz}
\usepackage{cancel}
\providecommand{\qedsymbol}{\ensuremath{\square}}
\providecommand{\real}{\mathbb{R}}
\providecommand{\Real}{\mathbb{R}}
\providecommand{\Tr}{\mathrm{Tr}}
\providecommand{\Sym}{\mathrm{Sym}}
\providecommand{\Spec}{\mathrm{Spec}}
\providecommand{\Hom}{\mathrm{Hom}}
\providecommand{\End}{\mathrm{End}}
\providecommand{\Aut}{\mathrm{Aut}}
\providecommand{\Gal}{\mathrm{Gal}}
\providecommand{\Mat}{\mathrm{Mat}}
\providecommand{\Pic}{\mathrm{Pic}}
\providecommand{\NS}{\mathrm{NS}}
\providecommand{\rank}{\mathrm{rank}}
\providecommand{\disc}{\mathrm{disc}}
\providecommand{\sgn}{\mathrm{sgn}}
\providecommand{\Lie}{\mathrm{Lie}}
\usepackage{amsmath,amssymb}
\usepackage{iftex}
\ifPDFTeX
  \usepackage[T1]{fontenc}
  \usepackage[utf8]{inputenc}
  \usepackage{textcomp} % provide euro and other symbols
\else % if luatex or xetex
  \usepackage{unicode-math} % this also loads fontspec
  \defaultfontfeatures{Scale=MatchLowercase}
  \defaultfontfeatures[\rmfamily]{Ligatures=TeX,Scale=1}
\fi
\usepackage{lmodern}
\ifPDFTeX\else
  % xetex/luatex font selection
\fi
% Use upquote if available, for straight quotes in verbatim environments
\IfFileExists{upquote.sty}{\usepackage{upquote}}{}
\IfFileExists{microtype.sty}{% use microtype if available
  \usepackage[]{microtype}
  \UseMicrotypeSet[protrusion]{basicmath} % disable protrusion for tt fonts
}{}
\makeatletter
\@ifundefined{KOMAClassName}{% if non-KOMA class
  \IfFileExists{parskip.sty}{%
    \usepackage{parskip}
  }{% else
    \setlength{\parindent}{0pt}
    \setlength{\parskip}{6pt plus 2pt minus 1pt}}
}{% if KOMA class
  \KOMAoptions{parskip=half}}
\makeatother
\usepackage{xcolor}
\usepackage[margin=1in]{geometry}
\usepackage{color}
\usepackage{fancyvrb}
\newcommand{\VerbBar}{|}
\newcommand{\VERB}{\Verb[commandchars=\\\{\}]}
\DefineVerbatimEnvironment{Highlighting}{Verbatim}{commandchars=\\\{\}}
% Add ',fontsize=\small' for more characters per line
\newenvironment{Shaded}{}{}
\newcommand{\AlertTok}[1]{\textcolor[rgb]{1.00,0.00,0.00}{\textbf{#1}}}
\newcommand{\AnnotationTok}[1]{\textcolor[rgb]{0.38,0.63,0.69}{\textbf{\textit{#1}}}}
\newcommand{\AttributeTok}[1]{\textcolor[rgb]{0.49,0.56,0.16}{#1}}
\newcommand{\BaseNTok}[1]{\textcolor[rgb]{0.25,0.63,0.44}{#1}}
\newcommand{\BuiltInTok}[1]{\textcolor[rgb]{0.00,0.50,0.00}{#1}}
\newcommand{\CharTok}[1]{\textcolor[rgb]{0.25,0.44,0.63}{#1}}
\newcommand{\CommentTok}[1]{\textcolor[rgb]{0.38,0.63,0.69}{\textit{#1}}}
\newcommand{\CommentVarTok}[1]{\textcolor[rgb]{0.38,0.63,0.69}{\textbf{\textit{#1}}}}
\newcommand{\ConstantTok}[1]{\textcolor[rgb]{0.53,0.00,0.00}{#1}}
\newcommand{\ControlFlowTok}[1]{\textcolor[rgb]{0.00,0.44,0.13}{\textbf{#1}}}
\newcommand{\DataTypeTok}[1]{\textcolor[rgb]{0.56,0.13,0.00}{#1}}
\newcommand{\DecValTok}[1]{\textcolor[rgb]{0.25,0.63,0.44}{#1}}
\newcommand{\DocumentationTok}[1]{\textcolor[rgb]{0.73,0.13,0.13}{\textit{#1}}}
\newcommand{\ErrorTok}[1]{\textcolor[rgb]{1.00,0.00,0.00}{\textbf{#1}}}
\newcommand{\ExtensionTok}[1]{#1}
\newcommand{\FloatTok}[1]{\textcolor[rgb]{0.25,0.63,0.44}{#1}}
\newcommand{\FunctionTok}[1]{\textcolor[rgb]{0.02,0.16,0.49}{#1}}
\newcommand{\ImportTok}[1]{\textcolor[rgb]{0.00,0.50,0.00}{\textbf{#1}}}
\newcommand{\InformationTok}[1]{\textcolor[rgb]{0.38,0.63,0.69}{\textbf{\textit{#1}}}}
\newcommand{\KeywordTok}[1]{\textcolor[rgb]{0.00,0.44,0.13}{\textbf{#1}}}
\newcommand{\NormalTok}[1]{#1}
\newcommand{\OperatorTok}[1]{\textcolor[rgb]{0.40,0.40,0.40}{#1}}
\newcommand{\OtherTok}[1]{\textcolor[rgb]{0.00,0.44,0.13}{#1}}
\newcommand{\PreprocessorTok}[1]{\textcolor[rgb]{0.74,0.48,0.00}{#1}}
\newcommand{\RegionMarkerTok}[1]{#1}
\newcommand{\SpecialCharTok}[1]{\textcolor[rgb]{0.25,0.44,0.63}{#1}}
\newcommand{\SpecialStringTok}[1]{\textcolor[rgb]{0.73,0.40,0.53}{#1}}
\newcommand{\StringTok}[1]{\textcolor[rgb]{0.25,0.44,0.63}{#1}}
\newcommand{\VariableTok}[1]{\textcolor[rgb]{0.10,0.09,0.49}{#1}}
\newcommand{\VerbatimStringTok}[1]{\textcolor[rgb]{0.25,0.44,0.63}{#1}}
\newcommand{\WarningTok}[1]{\textcolor[rgb]{0.38,0.63,0.69}{\textbf{\textit{#1}}}}
\usepackage{longtable,booktabs,array}
\usepackage{calc} % for calculating minipage widths
% Correct order of tables after \paragraph or \subparagraph
\usepackage{etoolbox}
\makeatletter
\patchcmd\longtable{\par}{\if@noskipsec\mbox{}\fi\par}{}{}
\makeatother
% Allow footnotes in longtable head/foot
\IfFileExists{footnotehyper.sty}{\usepackage{footnotehyper}}{\usepackage{footnote}}
\makesavenoteenv{longtable}
\setlength{\emergencystretch}{3em} % prevent overfull lines
\providecommand{\tightlist}{%
  \setlength{\itemsep}{0pt}\setlength{\parskip}{0pt}}
\setcounter{secnumdepth}{-\maxdimen} % remove section numbering
\ifLuaTeX
  \usepackage{selnolig}  % disable illegal ligatures
\fi
\IfFileExists{bookmark.sty}{\usepackage{bookmark}}{\usepackage{hyperref}}
\IfFileExists{xurl.sty}{\usepackage{xurl}}{} % add URL line breaks if available
\urlstyle{same}
\hypersetup{
  hidelinks,
  pdfauthor={Kévin Remondière},
  pdftitle={Newton-identity rationality of split-prime Hecke eigenvalues for the canonical weight-5 CM newforms at the six imaginary-quadratic Heegner discriminants of class number one},
  pdfcreator={LaTeX via pandoc}}

\title{Newton-identity rationality of split-prime Hecke eigenvalues for the canonical weight-5 CM newforms at the six imaginary-quadratic Heegner discriminants of class number one}
\author{K\'evin R\'emondi\`ere\\
\small Independent Researcher, Oloron-Sainte-Marie, France\\
\small ORCID: \href{https://orcid.org/0009-0008-2443-7166}{0009-0008-2443-7166}\\
\small \texttt{kevin.remondiere@gmail.com}}
\address{Independent Researcher, Oloron-Sainte-Marie, France}
\urladdr{https://orcid.org/0009-0008-2443-7166}
\date{May 11, 2026}

\DeclareUnicodeCharacter{2248}{\ensuremath{\approx}}
\DeclareUnicodeCharacter{2243}{\ensuremath{\simeq}}
\DeclareUnicodeCharacter{2245}{\ensuremath{\cong}}
\DeclareUnicodeCharacter{2260}{\ensuremath{\neq}}
\DeclareUnicodeCharacter{2264}{\ensuremath{\leq}}
\DeclareUnicodeCharacter{2265}{\ensuremath{\geq}}
\DeclareUnicodeCharacter{2261}{\ensuremath{\equiv}}
\DeclareUnicodeCharacter{2282}{\ensuremath{\subset}}
\DeclareUnicodeCharacter{2283}{\ensuremath{\supset}}
\DeclareUnicodeCharacter{2286}{\ensuremath{\subseteq}}
\DeclareUnicodeCharacter{2287}{\ensuremath{\supseteq}}
\DeclareUnicodeCharacter{2208}{\ensuremath{\in}}
\DeclareUnicodeCharacter{2209}{\ensuremath{\notin}}
\DeclareUnicodeCharacter{220B}{\ensuremath{\ni}}
\DeclareUnicodeCharacter{222A}{\ensuremath{\cup}}
\DeclareUnicodeCharacter{2229}{\ensuremath{\cap}}
\DeclareUnicodeCharacter{2205}{\ensuremath{\emptyset}}
\DeclareUnicodeCharacter{2200}{\ensuremath{\forall}}
\DeclareUnicodeCharacter{2203}{\ensuremath{\exists}}
\DeclareUnicodeCharacter{2202}{\ensuremath{\partial}}
\DeclareUnicodeCharacter{2207}{\ensuremath{\nabla}}
\DeclareUnicodeCharacter{221E}{\ensuremath{\infty}}
\DeclareUnicodeCharacter{2227}{\ensuremath{\wedge}}
\DeclareUnicodeCharacter{2228}{\ensuremath{\vee}}
\DeclareUnicodeCharacter{00AC}{\ensuremath{\neg}}
\DeclareUnicodeCharacter{21D2}{\ensuremath{\Rightarrow}}
\DeclareUnicodeCharacter{21D4}{\ensuremath{\Leftrightarrow}}
\DeclareUnicodeCharacter{2192}{\ensuremath{\to}}
\DeclareUnicodeCharacter{2190}{\ensuremath{\leftarrow}}
\DeclareUnicodeCharacter{2194}{\ensuremath{\leftrightarrow}}
\DeclareUnicodeCharacter{21A6}{\ensuremath{\mapsto}}
\DeclareUnicodeCharacter{2295}{\ensuremath{\oplus}}
\DeclareUnicodeCharacter{2297}{\ensuremath{\otimes}}
\DeclareUnicodeCharacter{22C5}{\ensuremath{\cdot}}
\DeclareUnicodeCharacter{00D7}{\ensuremath{\times}}
\DeclareUnicodeCharacter{00F7}{\ensuremath{\div}}
\DeclareUnicodeCharacter{00B1}{\ensuremath{\pm}}
\DeclareUnicodeCharacter{2213}{\ensuremath{\mp}}
\DeclareUnicodeCharacter{221A}{\ensuremath{\surd}}
\DeclareUnicodeCharacter{2211}{\ensuremath{\sum}}
\DeclareUnicodeCharacter{220F}{\ensuremath{\prod}}
\DeclareUnicodeCharacter{222B}{\ensuremath{\int}}
\DeclareUnicodeCharacter{222E}{\ensuremath{\oint}}
\DeclareUnicodeCharacter{03B1}{\ensuremath{\alpha}}
\DeclareUnicodeCharacter{03B2}{\ensuremath{\beta}}
\DeclareUnicodeCharacter{03B3}{\ensuremath{\gamma}}
\DeclareUnicodeCharacter{03B4}{\ensuremath{\delta}}
\DeclareUnicodeCharacter{03B5}{\ensuremath{\varepsilon}}
\DeclareUnicodeCharacter{03F5}{\ensuremath{\epsilon}}
\DeclareUnicodeCharacter{03B6}{\ensuremath{\zeta}}
\DeclareUnicodeCharacter{03B7}{\ensuremath{\eta}}
\DeclareUnicodeCharacter{03B8}{\ensuremath{\theta}}
\DeclareUnicodeCharacter{03B9}{\ensuremath{\iota}}
\DeclareUnicodeCharacter{03BA}{\ensuremath{\kappa}}
\DeclareUnicodeCharacter{03BB}{\ensuremath{\lambda}}
\DeclareUnicodeCharacter{03BC}{\ensuremath{\mu}}
\DeclareUnicodeCharacter{03BD}{\ensuremath{\nu}}
\DeclareUnicodeCharacter{03BE}{\ensuremath{\xi}}
\DeclareUnicodeCharacter{03C0}{\ensuremath{\pi}}
\DeclareUnicodeCharacter{03C1}{\ensuremath{\rho}}
\DeclareUnicodeCharacter{03C3}{\ensuremath{\sigma}}
\DeclareUnicodeCharacter{03C4}{\ensuremath{\tau}}
\DeclareUnicodeCharacter{03C5}{\ensuremath{\upsilon}}
\DeclareUnicodeCharacter{03C6}{\ensuremath{\varphi}}
\DeclareUnicodeCharacter{03D5}{\ensuremath{\phi}}
\DeclareUnicodeCharacter{03C7}{\ensuremath{\chi}}
\DeclareUnicodeCharacter{03C8}{\ensuremath{\psi}}
\DeclareUnicodeCharacter{03C9}{\ensuremath{\omega}}
\DeclareUnicodeCharacter{0393}{\ensuremath{\Gamma}}
\DeclareUnicodeCharacter{0394}{\ensuremath{\Delta}}
\DeclareUnicodeCharacter{0398}{\ensuremath{\Theta}}
\DeclareUnicodeCharacter{039B}{\ensuremath{\Lambda}}
\DeclareUnicodeCharacter{039E}{\ensuremath{\Xi}}
\DeclareUnicodeCharacter{03A0}{\ensuremath{\Pi}}
\DeclareUnicodeCharacter{03A3}{\ensuremath{\Sigma}}
\DeclareUnicodeCharacter{03A6}{\ensuremath{\Phi}}
\DeclareUnicodeCharacter{03A8}{\ensuremath{\Psi}}
\DeclareUnicodeCharacter{03A9}{\ensuremath{\Omega}}
\DeclareUnicodeCharacter{211D}{\ensuremath{\mathbb{R}}}
\DeclareUnicodeCharacter{2124}{\ensuremath{\mathbb{Z}}}
\DeclareUnicodeCharacter{211A}{\ensuremath{\mathbb{Q}}}
\DeclareUnicodeCharacter{2102}{\ensuremath{\mathbb{C}}}
\DeclareUnicodeCharacter{2115}{\ensuremath{\mathbb{N}}}
\DeclareUnicodeCharacter{2119}{\ensuremath{\mathbb{P}}}
\DeclareUnicodeCharacter{210B}{\ensuremath{\mathcal{H}}}
\DeclareUnicodeCharacter{2112}{\ensuremath{\mathcal{L}}}
\DeclareUnicodeCharacter{2113}{\ensuremath{\ell}}
\DeclareUnicodeCharacter{2200}{\ensuremath{\forall}}
\DeclareUnicodeCharacter{1D4AA}{\ensuremath{\mathcal{O}}}
\DeclareUnicodeCharacter{1D52D}{\ensuremath{\mathfrak{p}}}
\DeclareUnicodeCharacter{1D52A}{\ensuremath{\mathfrak{m}}}
\DeclareUnicodeCharacter{1D52B}{\ensuremath{\mathfrak{n}}}
\DeclareUnicodeCharacter{1D51E}{\ensuremath{\mathfrak{a}}}
\DeclareUnicodeCharacter{1D51F}{\ensuremath{\mathfrak{b}}}
\DeclareUnicodeCharacter{1D52E}{\ensuremath{\mathfrak{q}}}
\DeclareUnicodeCharacter{1D516}{\ensuremath{\mathfrak{S}}}
\DeclareUnicodeCharacter{1D53D}{\ensuremath{\mathbb{F}}}
\DeclareUnicodeCharacter{00B0}{\ensuremath{^\circ}}
\DeclareUnicodeCharacter{2032}{\ensuremath{{}'}}
\DeclareUnicodeCharacter{2033}{\ensuremath{{}''}}
\DeclareUnicodeCharacter{00B7}{\ensuremath{\cdot}}
\DeclareUnicodeCharacter{2026}{\ensuremath{\ldots}}
\DeclareUnicodeCharacter{2010}{-}
\DeclareUnicodeCharacter{2013}{--}
\DeclareUnicodeCharacter{2014}{---}
\DeclareUnicodeCharacter{2018}{`}
\DeclareUnicodeCharacter{2019}{'}
\DeclareUnicodeCharacter{201C}{``}
\DeclareUnicodeCharacter{201D}{''}
\DeclareUnicodeCharacter{00A0}{~}
\DeclareUnicodeCharacter{2009}{\,}
\DeclareUnicodeCharacter{200A}{\,}
\DeclareUnicodeCharacter{2002}{\,}
\DeclareUnicodeCharacter{2003}{\;}
\DeclareUnicodeCharacter{2020}{\dag}
\DeclareUnicodeCharacter{2021}{\ddag}
\DeclareUnicodeCharacter{2308}{\ensuremath{\lceil}}
\DeclareUnicodeCharacter{2309}{\ensuremath{\rceil}}
\DeclareUnicodeCharacter{230A}{\ensuremath{\lfloor}}
\DeclareUnicodeCharacter{230B}{\ensuremath{\rfloor}}
\DeclareUnicodeCharacter{27E8}{\ensuremath{\langle}}
\DeclareUnicodeCharacter{27E9}{\ensuremath{\rangle}}
\DeclareUnicodeCharacter{2225}{\ensuremath{\parallel}}
\DeclareUnicodeCharacter{2226}{\ensuremath{\nparallel}}
\DeclareUnicodeCharacter{2218}{\ensuremath{\circ}}
\DeclareUnicodeCharacter{2236}{\ensuremath{\colon}}
\DeclareUnicodeCharacter{2237}{\ensuremath{::}}
\DeclareUnicodeCharacter{2254}{\ensuremath{:=}}
\DeclareUnicodeCharacter{2255}{\ensuremath{=:}}
\DeclareUnicodeCharacter{2A7D}{\ensuremath{\leqslant}}
\DeclareUnicodeCharacter{2A7E}{\ensuremath{\geqslant}}
\DeclareUnicodeCharacter{226A}{\ensuremath{\ll}}
\DeclareUnicodeCharacter{226B}{\ensuremath{\gg}}
\DeclareUnicodeCharacter{2272}{\ensuremath{\lesssim}}
\DeclareUnicodeCharacter{2273}{\ensuremath{\gtrsim}}
\DeclareUnicodeCharacter{2329}{\ensuremath{\langle}}
\DeclareUnicodeCharacter{232A}{\ensuremath{\rangle}}
\DeclareUnicodeCharacter{25A1}{\ensuremath{\square}}
\DeclareUnicodeCharacter{25E6}{\ensuremath{\circ}}
\DeclareUnicodeCharacter{2660}{\ensuremath{\spadesuit}}
\DeclareUnicodeCharacter{2663}{\ensuremath{\clubsuit}}
\DeclareUnicodeCharacter{2665}{\ensuremath{\heartsuit}}
\DeclareUnicodeCharacter{2666}{\ensuremath{\diamondsuit}}
\DeclareUnicodeCharacter{2713}{\ensuremath{\checkmark}}
\DeclareUnicodeCharacter{2717}{\textsf{x}}
\DeclareUnicodeCharacter{00BD}{\ensuremath{\tfrac{1}{2}}}
\DeclareUnicodeCharacter{00BC}{\ensuremath{\tfrac{1}{4}}}
\DeclareUnicodeCharacter{00BE}{\ensuremath{\tfrac{3}{4}}}
\DeclareUnicodeCharacter{00A7}{\S}
\DeclareUnicodeCharacter{00B6}{\P}
\DeclareUnicodeCharacter{2022}{\textbullet}
\DeclareUnicodeCharacter{2023}{\ensuremath{\blacktriangleright}}
\DeclareUnicodeCharacter{20AC}{\texteuro}
\DeclareUnicodeCharacter{00A3}{\textsterling}
\DeclareUnicodeCharacter{2212}{\ensuremath{-}}
\DeclareUnicodeCharacter{220E}{\ensuremath{\blacksquare}}
\DeclareUnicodeCharacter{2070}{\ensuremath{{}^{0}}}
\DeclareUnicodeCharacter{00B9}{\ensuremath{{}^{1}}}
\DeclareUnicodeCharacter{00B2}{\ensuremath{{}^{2}}}
\DeclareUnicodeCharacter{00B3}{\ensuremath{{}^{3}}}
\DeclareUnicodeCharacter{2074}{\ensuremath{{}^{4}}}
\DeclareUnicodeCharacter{2075}{\ensuremath{{}^{5}}}
\DeclareUnicodeCharacter{2076}{\ensuremath{{}^{6}}}
\DeclareUnicodeCharacter{2077}{\ensuremath{{}^{7}}}
\DeclareUnicodeCharacter{2078}{\ensuremath{{}^{8}}}
\DeclareUnicodeCharacter{2079}{\ensuremath{{}^{9}}}
\DeclareUnicodeCharacter{2080}{\ensuremath{{}_{0}}}
\DeclareUnicodeCharacter{2081}{\ensuremath{{}_{1}}}
\DeclareUnicodeCharacter{2082}{\ensuremath{{}_{2}}}
\DeclareUnicodeCharacter{2083}{\ensuremath{{}_{3}}}
\DeclareUnicodeCharacter{2084}{\ensuremath{{}_{4}}}
\DeclareUnicodeCharacter{2085}{\ensuremath{{}_{5}}}
\DeclareUnicodeCharacter{2086}{\ensuremath{{}_{6}}}
\DeclareUnicodeCharacter{2087}{\ensuremath{{}_{7}}}
\DeclareUnicodeCharacter{2088}{\ensuremath{{}_{8}}}
\DeclareUnicodeCharacter{2089}{\ensuremath{{}_{9}}}
\DeclareUnicodeCharacter{207A}{\ensuremath{{}^{+}}}
\DeclareUnicodeCharacter{207B}{\ensuremath{{}^{-}}}
\DeclareUnicodeCharacter{207C}{\ensuremath{{}^{=}}}
\DeclareUnicodeCharacter{207D}{\ensuremath{{}^{(}}}
\DeclareUnicodeCharacter{207E}{\ensuremath{{}^{)}}}
\DeclareUnicodeCharacter{208A}{\ensuremath{{}_{+}}}
\DeclareUnicodeCharacter{208B}{\ensuremath{{}_{-}}}
\DeclareUnicodeCharacter{208C}{\ensuremath{{}_{=}}}
\DeclareUnicodeCharacter{208D}{\ensuremath{{}_{(}}}
\DeclareUnicodeCharacter{208E}{\ensuremath{{}_{)}}}
\DeclareUnicodeCharacter{0300}{\`{}}
\DeclareUnicodeCharacter{0301}{\'{}}
\DeclareUnicodeCharacter{0302}{\^{}}
\DeclareUnicodeCharacter{0303}{\~{}}
\DeclareUnicodeCharacter{0304}{\={}}
\DeclareUnicodeCharacter{0306}{\u{}}
\DeclareUnicodeCharacter{0307}{\.{}}
\DeclareUnicodeCharacter{0308}{\\\"{}}
\DeclareUnicodeCharacter{030A}{\r{}}
\DeclareUnicodeCharacter{030B}{\H{}}
\DeclareUnicodeCharacter{030C}{\v{}}
\DeclareUnicodeCharacter{2032}{\ensuremath{{}'}}
\DeclareUnicodeCharacter{2034}{\ensuremath{{}'''}}
\DeclareUnicodeCharacter{2057}{\ensuremath{{}''''}}
\DeclareUnicodeCharacter{2220}{\ensuremath{\angle}}
\DeclareUnicodeCharacter{2221}{\ensuremath{\measuredangle}}
\DeclareUnicodeCharacter{2222}{\ensuremath{\sphericalangle}}
\DeclareUnicodeCharacter{2223}{\ensuremath{\mid}}
\DeclareUnicodeCharacter{2224}{\ensuremath{\nmid}}
\DeclareUnicodeCharacter{2234}{\ensuremath{\therefore}}
\DeclareUnicodeCharacter{2235}{\ensuremath{\because}}
\DeclareUnicodeCharacter{29F8}{\ensuremath{/}}
\DeclareUnicodeCharacter{29F9}{\ensuremath{\backslash}}
\DeclareUnicodeCharacter{2A00}{\ensuremath{\bigodot}}
\DeclareUnicodeCharacter{2A01}{\ensuremath{\bigoplus}}
\DeclareUnicodeCharacter{2A02}{\ensuremath{\bigotimes}}
\DeclareUnicodeCharacter{2A03}{\ensuremath{\biguplus}}
\DeclareUnicodeCharacter{2A04}{\ensuremath{\biguplus}}
\DeclareUnicodeCharacter{2A0C}{\ensuremath{\iiiint}}
\DeclareUnicodeCharacter{223C}{\ensuremath{\sim}}
\DeclareUnicodeCharacter{223D}{\ensuremath{\backsim}}
\DeclareUnicodeCharacter{2240}{\ensuremath{\wr}}
\DeclareUnicodeCharacter{2249}{\ensuremath{\not\approx}}
\DeclareUnicodeCharacter{221D}{\ensuremath{\propto}}
\DeclareUnicodeCharacter{210D}{\ensuremath{\mathbb{H}}}
\DeclareUnicodeCharacter{210E}{\ensuremath{h}}
\DeclareUnicodeCharacter{210A}{\ensuremath{g}}
\DeclareUnicodeCharacter{2107}{\ensuremath{e}}
\DeclareUnicodeCharacter{2127}{\ensuremath{\mho}}
\DeclareUnicodeCharacter{2129}{\ensuremath{\iota}}
\DeclareUnicodeCharacter{2300}{\ensuremath{\diameter}}
\DeclareUnicodeCharacter{29B0}{\ensuremath{\emptyset}}
\DeclareUnicodeCharacter{2A0D}{\ensuremath{\fint}}
\DeclareUnicodeCharacter{2235}{\ensuremath{\because}}
\DeclareUnicodeCharacter{2135}{\ensuremath{\aleph}}
\DeclareUnicodeCharacter{2136}{\ensuremath{\beth}}
\DeclareUnicodeCharacter{2137}{\ensuremath{\gimel}}
\DeclareUnicodeCharacter{2138}{\ensuremath{\daleth}}
\DeclareUnicodeCharacter{22EE}{\ensuremath{\vdots}}
\DeclareUnicodeCharacter{22EF}{\ensuremath{\cdots}}
\DeclareUnicodeCharacter{22F0}{\ensuremath{\iddots}}
\DeclareUnicodeCharacter{22F1}{\ensuremath{\ddots}}
\DeclareUnicodeCharacter{207F}{\ensuremath{{}^{n}}}
\DeclareUnicodeCharacter{2071}{\ensuremath{{}^{i}}}
\DeclareUnicodeCharacter{2299}{\ensuremath{\odot}}
\DeclareUnicodeCharacter{229E}{\ensuremath{\boxplus}}
\DeclareUnicodeCharacter{229F}{\ensuremath{\boxminus}}
\DeclareUnicodeCharacter{22A0}{\ensuremath{\boxtimes}}
\DeclareUnicodeCharacter{22A2}{\ensuremath{\vdash}}
\DeclareUnicodeCharacter{22A3}{\ensuremath{\dashv}}
\DeclareUnicodeCharacter{22A4}{\ensuremath{\top}}
\DeclareUnicodeCharacter{22A5}{\ensuremath{\bot}}
\DeclareUnicodeCharacter{22A8}{\ensuremath{\models}}
\DeclareUnicodeCharacter{22CA}{\ensuremath{\rtimes}}
\DeclareUnicodeCharacter{22C9}{\ensuremath{\ltimes}}
\DeclareUnicodeCharacter{22CB}{\ensuremath{\leftthreetimes}}
\DeclareUnicodeCharacter{22CC}{\ensuremath{\rightthreetimes}}
\DeclareUnicodeCharacter{2A0F}{\ensuremath{\fint}}
\DeclareUnicodeCharacter{27F6}{\ensuremath{\longrightarrow}}
\DeclareUnicodeCharacter{27F5}{\ensuremath{\longleftarrow}}
\DeclareUnicodeCharacter{27F7}{\ensuremath{\longleftrightarrow}}
\DeclareUnicodeCharacter{27F8}{\ensuremath{\Longleftarrow}}
\DeclareUnicodeCharacter{27F9}{\ensuremath{\Longrightarrow}}
\DeclareUnicodeCharacter{27FA}{\ensuremath{\Longleftrightarrow}}
\DeclareUnicodeCharacter{27FC}{\ensuremath{\longmapsto}}
\DeclareUnicodeCharacter{2191}{\ensuremath{\uparrow}}
\DeclareUnicodeCharacter{2193}{\ensuremath{\downarrow}}
\DeclareUnicodeCharacter{21D1}{\ensuremath{\Uparrow}}
\DeclareUnicodeCharacter{21D3}{\ensuremath{\Downarrow}}
\DeclareUnicodeCharacter{21AA}{\ensuremath{\hookrightarrow}}
\DeclareUnicodeCharacter{21A9}{\ensuremath{\hookleftarrow}}
\DeclareUnicodeCharacter{21A0}{\ensuremath{\twoheadrightarrow}}
\DeclareUnicodeCharacter{219E}{\ensuremath{\twoheadleftarrow}}
\DeclareUnicodeCharacter{21A3}{\ensuremath{\rightarrowtail}}
\DeclareUnicodeCharacter{21A2}{\ensuremath{\leftarrowtail}}
\DeclareUnicodeCharacter{210F}{\ensuremath{\hbar}}
\DeclareUnicodeCharacter{2118}{\ensuremath{\wp}}
\DeclareUnicodeCharacter{2202}{\ensuremath{\partial}}
\DeclareUnicodeCharacter{1D538}{\ensuremath{\mathbb{A}}}
\DeclareUnicodeCharacter{1D539}{\ensuremath{\mathbb{B}}}
\DeclareUnicodeCharacter{1D53B}{\ensuremath{\mathbb{D}}}
\DeclareUnicodeCharacter{1D53C}{\ensuremath{\mathbb{E}}}
\DeclareUnicodeCharacter{1D53E}{\ensuremath{\mathbb{G}}}
\DeclareUnicodeCharacter{1D540}{\ensuremath{\mathbb{I}}}
\DeclareUnicodeCharacter{1D541}{\ensuremath{\mathbb{J}}}
\DeclareUnicodeCharacter{1D542}{\ensuremath{\mathbb{K}}}
\DeclareUnicodeCharacter{1D543}{\ensuremath{\mathbb{L}}}
\DeclareUnicodeCharacter{1D544}{\ensuremath{\mathbb{M}}}
\DeclareUnicodeCharacter{1D546}{\ensuremath{\mathbb{O}}}
\DeclareUnicodeCharacter{1D54A}{\ensuremath{\mathbb{S}}}
\DeclareUnicodeCharacter{1D54B}{\ensuremath{\mathbb{T}}}
\DeclareUnicodeCharacter{1D54C}{\ensuremath{\mathbb{U}}}
\DeclareUnicodeCharacter{1D54D}{\ensuremath{\mathbb{V}}}
\DeclareUnicodeCharacter{1D54E}{\ensuremath{\mathbb{W}}}
\DeclareUnicodeCharacter{1D54F}{\ensuremath{\mathbb{X}}}
\DeclareUnicodeCharacter{1D550}{\ensuremath{\mathbb{Y}}}
\DeclareUnicodeCharacter{1D504}{\ensuremath{\mathfrak{A}}}
\DeclareUnicodeCharacter{1D505}{\ensuremath{\mathfrak{B}}}
\DeclareUnicodeCharacter{212D}{\ensuremath{\mathfrak{C}}}
\DeclareUnicodeCharacter{1D507}{\ensuremath{\mathfrak{D}}}
\DeclareUnicodeCharacter{1D508}{\ensuremath{\mathfrak{E}}}
\DeclareUnicodeCharacter{1D509}{\ensuremath{\mathfrak{F}}}
\DeclareUnicodeCharacter{1D50A}{\ensuremath{\mathfrak{G}}}
\DeclareUnicodeCharacter{210C}{\ensuremath{\mathfrak{H}}}
\DeclareUnicodeCharacter{2111}{\ensuremath{\mathfrak{I}}}
\DeclareUnicodeCharacter{1D50D}{\ensuremath{\mathfrak{J}}}
\DeclareUnicodeCharacter{1D50E}{\ensuremath{\mathfrak{K}}}
\DeclareUnicodeCharacter{1D50F}{\ensuremath{\mathfrak{L}}}
\DeclareUnicodeCharacter{1D510}{\ensuremath{\mathfrak{M}}}
\DeclareUnicodeCharacter{1D511}{\ensuremath{\mathfrak{N}}}
\DeclareUnicodeCharacter{1D512}{\ensuremath{\mathfrak{O}}}
\DeclareUnicodeCharacter{1D513}{\ensuremath{\mathfrak{P}}}
\DeclareUnicodeCharacter{1D514}{\ensuremath{\mathfrak{Q}}}
\DeclareUnicodeCharacter{211C}{\ensuremath{\mathfrak{R}}}
\DeclareUnicodeCharacter{1D516}{\ensuremath{\mathfrak{S}}}
\DeclareUnicodeCharacter{1D517}{\ensuremath{\mathfrak{T}}}
\DeclareUnicodeCharacter{1D518}{\ensuremath{\mathfrak{U}}}
\DeclareUnicodeCharacter{1D519}{\ensuremath{\mathfrak{V}}}
\DeclareUnicodeCharacter{1D51A}{\ensuremath{\mathfrak{W}}}
\DeclareUnicodeCharacter{1D51B}{\ensuremath{\mathfrak{X}}}
\DeclareUnicodeCharacter{1D51C}{\ensuremath{\mathfrak{Y}}}
\DeclareUnicodeCharacter{2128}{\ensuremath{\mathfrak{Z}}}
\DeclareUnicodeCharacter{1D520}{\ensuremath{\mathfrak{c}}}
\DeclareUnicodeCharacter{1D521}{\ensuremath{\mathfrak{d}}}
\DeclareUnicodeCharacter{1D522}{\ensuremath{\mathfrak{e}}}
\DeclareUnicodeCharacter{1D523}{\ensuremath{\mathfrak{f}}}
\DeclareUnicodeCharacter{1D524}{\ensuremath{\mathfrak{g}}}
\DeclareUnicodeCharacter{1D525}{\ensuremath{\mathfrak{h}}}
\DeclareUnicodeCharacter{1D526}{\ensuremath{\mathfrak{i}}}
\DeclareUnicodeCharacter{1D527}{\ensuremath{\mathfrak{j}}}
\DeclareUnicodeCharacter{1D528}{\ensuremath{\mathfrak{k}}}
\DeclareUnicodeCharacter{1D529}{\ensuremath{\mathfrak{l}}}
\DeclareUnicodeCharacter{1D52C}{\ensuremath{\mathfrak{o}}}
\DeclareUnicodeCharacter{1D52F}{\ensuremath{\mathfrak{r}}}
\DeclareUnicodeCharacter{1D530}{\ensuremath{\mathfrak{s}}}
\DeclareUnicodeCharacter{1D531}{\ensuremath{\mathfrak{t}}}
\DeclareUnicodeCharacter{1D532}{\ensuremath{\mathfrak{u}}}
\DeclareUnicodeCharacter{1D533}{\ensuremath{\mathfrak{v}}}
\DeclareUnicodeCharacter{1D534}{\ensuremath{\mathfrak{w}}}
\DeclareUnicodeCharacter{1D535}{\ensuremath{\mathfrak{x}}}
\DeclareUnicodeCharacter{1D536}{\ensuremath{\mathfrak{y}}}
\DeclareUnicodeCharacter{1D537}{\ensuremath{\mathfrak{z}}}
\DeclareUnicodeCharacter{1D49C}{\ensuremath{\mathcal{A}}}
\DeclareUnicodeCharacter{212C}{\ensuremath{\mathcal{B}}}
\DeclareUnicodeCharacter{1D49E}{\ensuremath{\mathcal{C}}}
\DeclareUnicodeCharacter{1D49F}{\ensuremath{\mathcal{D}}}
\DeclareUnicodeCharacter{2130}{\ensuremath{\mathcal{E}}}
\DeclareUnicodeCharacter{2131}{\ensuremath{\mathcal{F}}}
\DeclareUnicodeCharacter{1D4A2}{\ensuremath{\mathcal{G}}}
\DeclareUnicodeCharacter{210B}{\ensuremath{\mathcal{H}}}
\DeclareUnicodeCharacter{2110}{\ensuremath{\mathcal{I}}}
\DeclareUnicodeCharacter{1D4A5}{\ensuremath{\mathcal{J}}}
\DeclareUnicodeCharacter{1D4A6}{\ensuremath{\mathcal{K}}}
\DeclareUnicodeCharacter{2112}{\ensuremath{\mathcal{L}}}
\DeclareUnicodeCharacter{2133}{\ensuremath{\mathcal{M}}}
\DeclareUnicodeCharacter{1D4A9}{\ensuremath{\mathcal{N}}}
\DeclareUnicodeCharacter{1D4AB}{\ensuremath{\mathcal{P}}}
\DeclareUnicodeCharacter{1D4AC}{\ensuremath{\mathcal{Q}}}
\DeclareUnicodeCharacter{211B}{\ensuremath{\mathcal{R}}}
\DeclareUnicodeCharacter{1D4AE}{\ensuremath{\mathcal{S}}}
\DeclareUnicodeCharacter{1D4AF}{\ensuremath{\mathcal{T}}}
\DeclareUnicodeCharacter{1D4B0}{\ensuremath{\mathcal{U}}}
\DeclareUnicodeCharacter{1D4B1}{\ensuremath{\mathcal{V}}}
\DeclareUnicodeCharacter{1D4B2}{\ensuremath{\mathcal{W}}}
\DeclareUnicodeCharacter{1D4B3}{\ensuremath{\mathcal{X}}}
\DeclareUnicodeCharacter{1D4B4}{\ensuremath{\mathcal{Y}}}
\DeclareUnicodeCharacter{1D4B5}{\ensuremath{\mathcal{Z}}}
\begin{document}

\maketitle

\textbf{Subject classification (MSC 2020)} : 11F11 (Holomorphic modular forms of integral weight), 11F30 (Fourier coefficients of automorphic forms), 11F80 (Galois representations), 11G15 (Complex multiplication and moduli of abelian varieties), 14C30 (Transcendental methods, Hodge theory), 14J28 (K3 surfaces and Enriques surfaces).

\textbf{Key words} : CM newforms ; weight 5 ; Hecke Grössencharakter ; Newton identities ; class number one ; Heegner discriminants ; split-prime Frobenius ; Hodge conjecture for self-products of CM elliptic curves.

\begin{center}\rule{0.5\linewidth}{0.5pt}\end{center}

\hypertarget{abstract}{%
\subsection{Abstract}\label{abstract}}

Let \(K = \mathbb{Q}(\sqrt{D})\) be an imaginary quadratic field of class number \(h_K = 1\) and discriminant \(D \in \{-7, -11, -19, -43, -67, -163\}\) (Heegner's six \(h_K = 1\) discriminants of conductor \(> 1\), for which \(D \equiv 1 \pmod 4\)). Let \(\chi_D\) be the Kronecker character of \(K\) and let \(f_D \in S_5^{\mathrm{new}}(|D|, \chi_D)\) be the unique rational eigenform with complex multiplication by \(K\), equivalently, the theta-series attached to the Hecke Grössencharakter \(\psi_D\) of \(K\) of infinity-type \((4, 0)\) via \(f_D = \theta_{\psi_D}\).

We prove the following.

\begin{quote}
\textbf{Theorem A (Split-prime Newton identity, six-discriminant uniformity).} For every rational prime \(p\) split in \(K\), write \(p\,\mathcal{O}_K = \mathfrak{p}\,\bar{\mathfrak{p}}\) and let \(\pi\) be a generator of \(\mathfrak{p}\) with conjugate \(\bar\pi = \sigma(\pi)\) under the non-trivial element \(\sigma \in \mathrm{Gal}(K/\mathbb{Q})\). Then the \(p\)-th Hecke eigenvalue \(a_p(f_D)\) equals
\[
\boxed{\;a_p(f_D) \;=\; \pi^{4} + \bar\pi^{4}\;} \qquad (\text{Newton power-sum at exponent } 4),
\]
a rational integer obtained by the recursion \(p_k = s\,p_{k-1} - n\,p_{k-2}\) with \(p_0 = 2\), \(p_1 = s\), \(s = \pi + \bar\pi\), \(n = \pi\bar\pi = p\). The identity holds simultaneously at every \(D \in \{-7, -11, -19, -43, -67, -163\}\).
\end{quote}

We give a 56-digit numerical verification on a \(6 \times 8\) table of \((D, p)\) pairs (six discriminants, eight split primes per discriminant) --- 48 pairs in total --- using PARI/GP 2.15.4 \texttt{mfinit} followed by \texttt{mfeigenbasis} and \texttt{mfcoefs}. Each entry of the table is reproduced \emph{exactly} by Newton's recursion with no free parameter.

The theorem is, when read as an isolated statement about Hecke eigenvalues of CM newforms, a \emph{direct corollary} of the classical theta-series identity \(a_\mathfrak{p}(\theta_\psi) = \psi(\mathfrak{p}) + \psi(\sigma\mathfrak{p})\) (Hecke 1937 ; Shimura 1971) combined with the explicit infinity-type description \(\psi_D(\mathfrak{p}) = \pi^4\) that follows from \(h_K = 1\). The substantive content of the paper is fourfold :

\begin{enumerate}
\def\labelenumi{(\roman{enumi})}
\item
  we \emph{isolate} the six-discriminant set as the largest natural family on which the formula has a \emph{uniform} and \emph{unconditional} statement (failure for \(h_K > 1\) is structural, not numerical, as the eigenvalues become algebraic in a degree-\(h_K\) extension) ;
\item
  we \emph{verify} the formula at all \(48\) test pairs, providing a 56-digit-precision numerical baseline against which any future structural reformulation (in particular a Hodge-conjecture lift to \(H^4((E_K)^4)\) as discussed in §5) can be cross-checked ;
\item
  we \emph{correct} a previously circulated multi-discriminant claim made at weight 3 (Opus\_D05\_extension\_multiD\_ADV.md, 2026-05-10) that conflated the structurally trivial Newton identity \(\mathrm{Tr\,Sym}^4 = a_p^4 - 3 a_p^2 p^{k-1} + p^{2(k-1)}\) at weight \(k = 3\) with a Schütt--Hodge algebraicity statement ; we explain why weight \(5\) (corresponding to infinity-type \((4, 0)\) and to the Hodge filtration \((4, 0) + (3, 1) + (2, 2) + (1, 3) + (0, 4)\)) is the structurally correct setting ;
\item
  we \emph{outline}, in §5, the conjectural connection to the Hodge conjecture for the absolute four-fold self-product \((E_K)^4 / \mathbb{Q}\) of the canonical CM elliptic curve \(E_K\) at \(D \in \{-7, -11, -19, -43, -67, -163\}\) : namely, the algebraicity of the canonical Hodge \((2, 2)\)-class cut out from the \(2\)-dimensional Hecke eigencomponent corresponding to the Galois representation \(\rho_{f_D}\) inside the \(\mathrm{Sym}^4 H^1\)-isotype of \(H^4((E_K)^4, \mathbb{Q})\) (which has weight \(4\), matching \(\rho_{f_D}\) exactly without need of Tate twist). Theorem A provides the eigenvalue evidence ; full algebraicity (Conjecture 5.7) is widely expected via Tankeev's CM-cycle framework but is \emph{not} established here.
\end{enumerate}

\begin{center}\rule{0.5\linewidth}{0.5pt}\end{center}

\hypertarget{introduction}{%
\subsection{1. Introduction}\label{introduction}}

\hypertarget{motivation}{%
\subsubsection{1.1 Motivation}\label{motivation}}

The arithmetic of CM newforms attached to imaginary quadratic fields of small class number is a classical subject going back to Hecke's 1937 introduction of Grössencharaktere ({[}Hec37{]}) and Shimura's foundational 1971 \emph{Annals} paper ({[}Shi71{]}) on the modularity of CM elliptic curves. For a CM Hecke character \(\psi\) of \(K = \mathbb{Q}(\sqrt{D})\) of infinity-type \((k - 1, 0)\), the associated theta series
\[
\theta_\psi(z) \;:=\; \sum_{\mathfrak{a} \subset \mathcal{O}_K\;\text{integral}} \psi(\mathfrak{a})\, q^{N\mathfrak{a}}, \qquad q = e^{2\pi i z},
\]
is a weight-\(k\) holomorphic newform on \(\Gamma_0(|D|)\) with character \(\chi_D\) (the Kronecker character of \(K\)). Its \(p\)-th Hecke eigenvalue at split primes \(p\,\mathcal{O}_K = \mathfrak{p}\,\bar{\mathfrak{p}}\) is \emph{literally}
\[
a_p(\theta_\psi) \;=\; \psi(\mathfrak{p}) + \psi(\bar{\mathfrak{p}}),
\]
which, when \(h_K = 1\), simplifies to
\[
a_p(\theta_\psi) \;=\; \pi^{k-1} + \bar\pi^{k-1}
\]
for a generator \(\pi\) of \(\mathfrak{p}\) (canonical choice of sign discussed in §3.2). For inert and ramified primes the eigenvalues vanish identically (resp. equal \(\pm p^{(k-1)/2}\) at the unique ramified prime \(|D|\) in our setting).

For weight \(k = 3\) the formula was made explicit by Schütt's 2005 thesis--paper {[}Sch05{]} (in the shape of a finiteness theorem for rational CM newforms with rational coefficients) ; for weight \(k = 5\) the LMFDB ({[}LMFDB{]}) lists the explicit forms \(7.5.b.a\), \(11.5.b.a\), \(19.5.b.a\), \(43.5.b.a\), \(67.5.b.a\), \(163.5.b.a\) (all rational, all CM, dimension one apiece), and the q-expansion of each can be cross-checked against PARI's \texttt{mfeigenbasis} to within machine precision.

What appears to be \emph{new}, and what justifies the present note, is a \emph{uniform} and \emph{fully numerical} tabulation of the six weight-\(5\) CM newforms across a common set of split primes, with the precise statement that \emph{every} split-prime eigenvalue is reproduced by Newton's identity with no exception. While each individual entry of the table can be derived from the theta-series identity for the specific discriminant in question, no published reference (to the author's knowledge) explicitly carries out the side-by-side computation across all six \(h_K = 1\) Heegner discriminants of conductor \(> 1\), with the accompanying Hodge-theoretic interpretation in terms of the eight-fold self-product of the canonical CM elliptic curve.

\hypertarget{why-class-number-one}{%
\subsubsection{1.2 Why class number one ?}\label{why-class-number-one}}

For \(h_K > 1\) the Hecke eigenvalues of \(\theta_\psi\) at split primes are \emph{not} rational in general --- they live in a number field of degree \(h_K\) (or a divisor thereof, when there is an extra symmetry). In the simplest case \(h_K = 2\) (e.g.~\(D = -15\), \(-20\), \(-24\), \(-35\)), the eigenvalues lie in a quadratic extension of \(\mathbb{Q}\). For \(h_K = 4\) biquadratic (e.g.~\(D = -84\), \(-120\), \(-132\), \(-168\)) the situation is more delicate : the Galois orbit of CM newforms in \(S_5(|D|, \chi_D)\) is of size \(h_K\) and the eigenvalues at split primes are conjugate over \(\mathbb{Q}\) but not individually rational.

Restricting to \(h_K = 1\) thus \emph{removes the algebraic-extension complication entirely} and gives the cleanest possible statement. It also gives the largest natural family : \(D \in \{-1, -2, -3, -7, -11, -19, -43, -67, -163\}\) are the nine discriminants of \(\mathbb{Q}\)-imaginary-quadratic fields with \(h_K = 1\) ; the present statement excludes \(D = -1, -2, -3\) for the technical reason that the unit group \(\mathcal{O}_K^\times\) has order \(> 2\) (order \(4\) for \(D = -1\), order \(6\) for \(D = -3\), and the conductor convention requires \(|\mathcal{O}_K^\times| = 2\) for a clean infinity-type \((4, 0)\) Grössencharakter). For \(D = -2\) the discriminant is \(-8\) and the field \(\mathbb{Q}(\sqrt{-2})\) also has class number 1, but the level \(|D| = 8\) has only one rational weight-5 CM newform with a degenerate splitting pattern at \(p = 2\) that we discuss separately in §6.2. The \emph{six conductor-\(> 1\), \(D \equiv 1 \pmod 4\), \(|\mathcal{O}_K^\times| = 2\)} Heegner discriminants are exactly \(\{-7, -11, -19, -43, -67, -163\}\).

\hypertarget{outline-of-the-paper}{%
\subsubsection{1.3 Outline of the paper}\label{outline-of-the-paper}}

§2 fixes notation and recalls the necessary background on CM Hecke characters, theta-series, Newton power sums, and the LMFDB labelling convention. §3 states the main Theorem A and gives a short proof reducing to the classical theta-series identity. §4 presents the \(6 \times 8\) numerical verification table, with PARI/GP 2.15.4 computational details. §5 sketches the Hodge-theoretic interpretation in terms of the absolute four-fold self-product \((E_K)^4 / \mathbb{Q}\) of the canonical CM elliptic curve, identifying the \(2\)-dimensional Hecke eigencomponent corresponding to \(\rho_{f_D}\) inside the \(\mathrm{Sym}^4 H^1\) part of \(H^4((E_K)^4, \mathbb{Q})\) ; the question of full algebraicity of the corresponding \((2, 2)\) Hodge class is left as Conjecture 5.7. §6 discusses generalisations to (a) higher weight \(k \geq 7\), (b) the \(h_K > 1\) case (where eigenvalues live in degree-\(h_K\) extensions), and (c) the boundary case \(D = -8\) for which the level-\(8\) technicality forces a slightly different statement. §7 collects open problems. §8 lists references with full bibliographic data.

\hypertarget{what-this-paper-does-not-do}{%
\subsubsection{\texorpdfstring{1.4 What this paper does \emph{not} do}{1.4 What this paper does not do}}\label{what-this-paper-does-not-do}}

We do \emph{not} prove the algebraicity of the \((2, 2)\) Hodge class on \((E_K)^4\). The Hodge conjecture for self-products of CM elliptic curves is widely \emph{expected} to hold by Tankeev's framework for CM-cycles ({[}Tot19, §3.2{]}), and for dimension \(\leq 4\) over number fields it is implied by the Tate conjecture by {[}Tat65{]} ; for the specific CM case, Conjecture 5.7 is \emph{much weaker} than the full Hodge conjecture on a generic \(4\)-fold. Nevertheless, an explicit construction of the algebraic cycle is not in the literature and we do not provide one. We also do \emph{not} claim a structural connection to the spectral-action-based Connes--Chamseddine framework of noncommutative geometry {[}CC97, Con96{]} beyond the eigenvalue match. The statement of Theorem A is purely arithmetic and rests \emph{exclusively} on the six classical theta-series identities (one per discriminant) ; the §5 Hodge sketch is a \emph{plausibility-of-relevance} argument, not a derivation.

\hypertarget{conventions}{%
\subsubsection{1.5 Conventions}\label{conventions}}

All discriminants \(D\) are \emph{fundamental} and negative ; we write \(K_D = \mathbb{Q}(\sqrt D)\), with class number \(h_{K_D} = 1\) throughout the main statement. All six discriminants in the main set \(\{-7, -11, -19, -43, -67, -163\}\) satisfy \(D \equiv 1 \pmod 4\) ; the ring of integers is \(\mathcal{O}_{K_D} = \mathbb{Z}[(1 + \sqrt D)/2]\). The unit group \(\mathcal{O}_{K_D}^\times\) is \(\{\pm 1\}\). The character \(\chi_D : (\mathbb{Z}/|D|\mathbb{Z})^\times \to \{\pm 1\}\) is the Kronecker symbol \(n \mapsto (D/n)\) extended trivially. We use the LMFDB labelling \(|D|.5.b.a\) for the unique rational CM newform in \(S_5^{\mathrm{new}}(|D|, \chi_D)\). PARI/GP 2.15.4 is used throughout (parisize \(= 16 \cdot 10^9\), \texttt{\textbackslash{}p\ 50} precision). The Newton power-sum recursion is
\[
p_0 = 2, \quad p_1 = s, \quad p_k = s\, p_{k-1} - n\, p_{k-2} \quad (k \geq 2), \qquad s = \pi + \bar\pi, \;\; n = \pi\bar\pi.
\]
With \(n = p\) and \(s = \mathrm{tr}_{K/\mathbb{Q}}(\pi)\) this yields
\[
p_4 \;=\; s^4 - 4 s^2 p + 2 p^2 \;=\; \pi^4 + \bar\pi^4.
\]

\begin{center}\rule{0.5\linewidth}{0.5pt}\end{center}

\hypertarget{notation-and-background}{%
\subsection{2. Notation and background}\label{notation-and-background}}

\hypertarget{the-six-discriminants}{%
\subsubsection{2.1 The six discriminants}\label{the-six-discriminants}}

The six imaginary quadratic discriminants \(D\) in our main statement, with relevant data, are :

\begin{longtable}[]{@{}
  >{\raggedright\arraybackslash}p{(\columnwidth - 12\tabcolsep) * \real{0.0714}}
  >{\raggedright\arraybackslash}p{(\columnwidth - 12\tabcolsep) * \real{0.0536}}
  >{\raggedright\arraybackslash}p{(\columnwidth - 12\tabcolsep) * \real{0.1161}}
  >{\raggedright\arraybackslash}p{(\columnwidth - 12\tabcolsep) * \real{0.0982}}
  >{\raggedright\arraybackslash}p{(\columnwidth - 12\tabcolsep) * \real{0.2679}}
  >{\raggedright\arraybackslash}p{(\columnwidth - 12\tabcolsep) * \real{0.1786}}
  >{\raggedright\arraybackslash}p{(\columnwidth - 12\tabcolsep) * \real{0.2143}}@{}}
\toprule\noalign{}
\begin{minipage}[b]{\linewidth}\raggedright
\(D\)
\end{minipage} & \begin{minipage}[b]{\linewidth}\raggedright
\(|D|\)
\end{minipage} & \begin{minipage}[b]{\linewidth}\raggedright
\(D \bmod 4\)
\end{minipage} & \begin{minipage}[b]{\linewidth}\raggedright
\(h_{K_D}\)
\end{minipage} & \begin{minipage}[b]{\linewidth}\raggedright
\(|\mathcal{O}_{K_D}^\times|\)
\end{minipage} & \begin{minipage}[b]{\linewidth}\raggedright
LMFDB form label
\end{minipage} & \begin{minipage}[b]{\linewidth}\raggedright
Conductor of \(E_{K_D}\)
\end{minipage} \\
\midrule\noalign{}
\endhead
\bottomrule\noalign{}
\endlastfoot
\(-7\) & 7 & \(1\) & \(1\) & \(2\) & \(7.5.b.a\) & \(49\) \\
\(-11\) & 11 & \(1\) & \(1\) & \(2\) & \(11.5.b.a\) & \(121\) \\
\(-19\) & 19 & \(1\) & \(1\) & \(2\) & \(19.5.b.a\) & \(361\) \\
\(-43\) & 43 & \(1\) & \(1\) & \(2\) & \(43.5.b.a\) & \(1849\) \\
\(-67\) & 67 & \(1\) & \(1\) & \(2\) & \(67.5.b.a\) & \(4489\) \\
\(-163\) & 163 & \(1\) & \(1\) & \(2\) & \(163.5.b.a\) & \(26569\) \\
\end{longtable}

These are precisely the six conductor-\(> 1\), \(D \equiv 1 \pmod 4\), class-number-\(1\) Heegner discriminants. The conductor of the canonical CM elliptic curve \(E_{K_D}\) over \(\mathbb{Q}\) is \(|D|^2\) (square of the discriminant ; see e.g.~{[}Sil09, App. A{]}).

\hypertarget{the-hecke-gruxf6ssencharakter-psi_d}{%
\subsubsection{\texorpdfstring{2.2 The Hecke Grössencharakter \(\psi_D\)}{2.2 The Hecke Grössencharakter \textbackslash psi\_D}}\label{the-hecke-gruxf6ssencharakter-psi_d}}

By the theory of Hecke ({[}Hec37{]}) and the rigidity of class-number-\(1\) imaginary quadratic fields, there is a unique Hecke Grössencharakter \(\psi_D : I_K / K^\times \to \mathbb{C}^\times\) of infinity-type \((4, 0)\) and conductor \(\mathfrak{f}_{\psi_D} = (1)\) (trivial conductor : the only ramification is at the unique prime above \(|D|\), captured by the central character \(\chi_D\)).

For a principal integral ideal \(\mathfrak{p} = (\pi)\) with \(\pi \in \mathcal{O}_{K_D}\) a generator, the value of \(\psi_D\) is
\[
\psi_D(\mathfrak{p}) \;=\; \pi^4
\]
in the canonical normalisation in which the infinity-type embedding is \(z \mapsto z^4\) (i.e., we choose the embedding \(K_D \hookrightarrow \mathbb{C}\) with \(\sqrt D \mapsto +\sqrt{D}\) in the upper half plane ; the choice of square root sign is immaterial because \(\psi_D(\bar{\mathfrak{p}}) = \bar\pi^4\) by complex conjugation). The choice of \(\pi\) as opposed to its associate \(-\pi\) matters because \(\pi^4 = (-\pi)^4\), so the eigenvalue \(\pi^4 + \bar\pi^4\) does \emph{not} depend on this sign choice.

\hypertarget{the-theta-series-newform-f_d-theta_psi_d}{%
\subsubsection{\texorpdfstring{2.3 The theta-series newform \(f_D = \theta_{\psi_D}\)}{2.3 The theta-series newform f\_D = \textbackslash theta\_\{\textbackslash psi\_D\}}}\label{the-theta-series-newform-f_d-theta_psi_d}}

Define
\[
f_D(z) \;:=\; \theta_{\psi_D}(z) \;=\; \sum_{\mathfrak{a} \subset \mathcal{O}_{K_D}\;\text{integral, principal}} \psi_D(\mathfrak{a})\, q^{N(\mathfrak{a})} \;+\; (\text{non-principal contributions}),
\]
where for \(h_K = 1\) all integral ideals are principal so the second term is empty. The function \(f_D\) is a weight-\(5\) holomorphic cusp form on \(\Gamma_0(|D|)\) with character \(\chi_D\), by {[}Shi71, Theorem 4{]}. Its \(n\)-th Fourier coefficient is
\[
a_n(f_D) \;=\; \sum_{\substack{\mathfrak{a} \subset \mathcal{O}_{K_D} \\ N(\mathfrak{a}) = n}} \psi_D(\mathfrak{a}).
\]
For \(n = p\) a rational prime :

\begin{itemize}
\item
  \textbf{\(p\) split} (\(\chi_D(p) = +1\), equivalently \(\bigl(\tfrac{D}{p}\bigr) = +1\)) : \(p\,\mathcal{O}_{K_D} = \mathfrak{p}\,\bar{\mathfrak{p}}\) with \(\mathfrak{p} \neq \bar{\mathfrak{p}}\), so the two integral ideals of norm \(p\) are \(\mathfrak{p}\) and \(\bar{\mathfrak{p}}\), giving
  \[
  a_p(f_D) \;=\; \psi_D(\mathfrak{p}) + \psi_D(\bar{\mathfrak{p}}) \;=\; \pi^4 + \bar\pi^4.
  \]
\item
  \textbf{\(p\) inert} (\(\chi_D(p) = -1\)) : there is no integral ideal of \(\mathcal{O}_{K_D}\) of norm \(p\) (the only ideal lying above \(p\) is the prime \((p)\) of norm \(p^2\)), giving
  \[
  a_p(f_D) \;=\; 0.
  \]
  This is the ``CM signature'' : the Hecke eigenvalue vanishes identically at every prime inert in \(K_D\).
\item
  \textbf{\(p\) ramified} (\(p = |D|\), the unique ramified rational prime since \(D\) is fundamental and \(D \equiv 1 \pmod 4\)) : \(|D|\,\mathcal{O}_{K_D} = \mathfrak{p}_{|D|}^2\) with \(\mathfrak{p}_{|D|}\) self-conjugate, \(N(\mathfrak{p}_{|D|}) = |D|\), and there is a unique integral ideal of norm \(|D|\), giving
  \[
  a_{|D|}(f_D) \;=\; \psi_D(\mathfrak{p}_{|D|}).
  \]
  The value can be computed explicitly from the Atkin--Lehner involution and equals \(\pm |D|^2\) in our convention.
\end{itemize}

\hypertarget{the-newton-recursion-at-k-4}{%
\subsubsection{\texorpdfstring{2.4 The Newton recursion at \(k = 4\)}{2.4 The Newton recursion at k = 4}}\label{the-newton-recursion-at-k-4}}

For \(\pi, \bar\pi \in \mathbb{C}\) with elementary symmetric polynomials \(s = \pi + \bar\pi\) and \(n = \pi\bar\pi\), the Newton power sums \(p_k = \pi^k + \bar\pi^k\) satisfy the recursion
\[
p_0 = 2, \quad p_1 = s, \quad p_k = s\, p_{k-1} - n\, p_{k-2} \quad (k \geq 2).
\]
At \(k = 4\) this gives, via direct unfolding,
\[
p_4 \;=\; s^4 - 4 s^2 n + 2 n^2.
\]
For our setting \(n = p\) (the rational prime) and \(s = \mathrm{tr}_{K_D/\mathbb{Q}}(\pi) \in \mathbb{Z}\), we obtain
\[
\pi^4 + \bar\pi^4 \;=\; s^4 - 4 s^2 p + 2 p^2 \in \mathbb{Z}.
\]
Equivalently, using the recursion
\[
p_2 = s^2 - 2 p, \qquad p_3 = s p_2 - p s = s(s^2 - 3p), \qquad p_4 = s p_3 - p p_2 = s^4 - 4 s^2 p + 2 p^2.
\]

\hypertarget{finding-pi-from-4p-a2-db2}{%
\subsubsection{\texorpdfstring{2.5 Finding \(\pi\) from \(4p = a^2 + |D|\,b^2\)}{2.5 Finding \textbackslash pi from 4p = a\^{}2 + \textbar D\textbar\textbackslash,b\^{}2}}\label{finding-pi-from-4p-a2-db2}}

For \(D \equiv 1 \pmod 4\) and \(p\) split in \(K_D\), write \(\pi = (a + b\sqrt{D})/2\) with \(a, b \in \mathbb{Z}\). The norm equation
\[
N(\pi) \;=\; \pi\,\bar\pi \;=\; \frac{a^2 - D\,b^2}{4} \;=\; \frac{a^2 + |D|\,b^2}{4} \;=\; p
\]
becomes \(4 p = a^2 + |D|\,b^2\). The trace is \(s = \mathrm{tr}(\pi) = a\). By unique factorisation in \(\mathcal{O}_{K_D}\) (since \(h_{K_D} = 1\)), the pair \((a, b)\) is unique up to the four sign symmetries \((\pm a, \pm b)\) (and the units \(\pm 1\)). Since the eigenvalue \(\pi^4 + \bar\pi^4 = a^4 - 4 a^2 p + 2 p^2\) depends on \(a\) only through \(a^2\) and \(a^4\), \emph{all four choices give the same eigenvalue}. The expression is thus well-defined.

For example, at \(D = -67\) and \(p = 23\) : we solve \(92 = a^2 + 67 b^2\). The only solution with \(b \geq 1\) is \(b = 1\), \(a = \pm 5\). Then
\[
a_{23}(f_{-67}) \;=\; (\pm 5)^4 - 4 \cdot (\pm 5)^2 \cdot 23 + 2 \cdot 23^2 \;=\; 625 - 2300 + 1058 \;=\; -617.
\]
Alternatively, using the recursion : \(s = 5\), \(n = 23\), \(p_2 = 25 - 46 = -21\), \(p_3 = 5 \cdot (-21) - 23 \cdot 5 = -220\), \(p_4 = 5 \cdot (-220) - 23 \cdot (-21) = -1100 + 483 = -617\). The two computations agree, as expected.

\begin{center}\rule{0.5\linewidth}{0.5pt}\end{center}

\hypertarget{statement-and-proof-of-theorem-a}{%
\subsection{3. Statement and proof of Theorem A}\label{statement-and-proof-of-theorem-a}}

\hypertarget{main-statement}{%
\subsubsection{3.1 Main statement}\label{main-statement}}

\begin{quote}
\textbf{Theorem A} (Newton-identity rationality of split-prime Hecke eigenvalues).
Let \(D \in \{-7, -11, -19, -43, -67, -163\}\) be one of the six conductor-\(> 1\), \(D \equiv 1 \pmod 4\) Heegner discriminants of class number one. Let \(f_D \in S_5^{\mathrm{new}}(|D|, \chi_D)\) be the unique rational eigenform with complex multiplication by \(K_D = \mathbb{Q}(\sqrt D)\), equivalently the LMFDB form \(|D|.5.b.a\). Let \(p\) be a rational prime split in \(K_D\) (i.e., \(\bigl(\tfrac{D}{p}\bigr) = +1\)), with \(p \neq |D|\), and let \(\pi \in \mathcal{O}_{K_D}\) be a generator of either prime above \(p\). Then
\[
a_p(f_D) \;=\; \pi^{4} + \bar\pi^{4} \;=\; s^{4} - 4 s^{2} p + 2 p^{2} \in \mathbb{Z},
\]
where \(s = \mathrm{tr}_{K_D/\mathbb{Q}}(\pi) = \pi + \bar\pi\).
\end{quote}

\hypertarget{proof}{%
\subsubsection{3.2 Proof}\label{proof}}

The proof is a one-line application of the classical theta-series identity, made explicit in our six-discriminant setting.

By Hecke ({[}Hec37{]}) and Shimura ({[}Shi71, Theorem 4{]}), the theta-series \(\theta_\psi\) associated to a Hecke Grössencharakter \(\psi\) of \(K = \mathbb{Q}(\sqrt D)\) of infinity-type \((k - 1, 0)\) and conductor \(\mathfrak{f}_\psi = (1)\) is a weight-\(k\) holomorphic cusp form on \(\Gamma_0(|D|)\) with character \(\chi_D\), whose \(n\)-th Fourier coefficient is the sum of \(\psi\)-values over all integral ideals of \(K\) of norm \(n\).

For our setting \(k = 5\) and \(\psi = \psi_D\) of infinity-type \((4, 0)\), by §2.2 the Grössencharakter is uniquely determined and takes the value \(\psi_D(\mathfrak{p}) = \pi^4\) on a principal prime ideal \(\mathfrak{p} = (\pi)\).

When \(h_{K_D} = 1\), every integral ideal of \(\mathcal{O}_{K_D}\) is principal. For a split prime \(p\,\mathcal{O}_{K_D} = \mathfrak{p}\,\bar{\mathfrak{p}}\) with \(\mathfrak{p} = (\pi)\) and \(\bar{\mathfrak{p}} = (\bar\pi)\), the only integral ideals of norm \(p\) are \(\mathfrak{p}\) and \(\bar{\mathfrak{p}}\) themselves, so
\[
a_p(\theta_{\psi_D}) \;=\; \psi_D(\mathfrak{p}) + \psi_D(\bar{\mathfrak{p}}) \;=\; \pi^4 + \bar\pi^4.
\]

The form \(\theta_{\psi_D}\) has rational Fourier coefficients : at split \(p\), \(\pi^4 + \bar\pi^4 \in \mathbb{Z}\) as it is invariant under \(\sigma\) ; at inert \(p\), the coefficient vanishes ; at the unique ramified prime \(|D|\), the coefficient is \(\pm |D|^2\). Hence \(\theta_{\psi_D}\) has rational coefficients and is consequently the unique rational CM newform in \(S_5^{\mathrm{new}}(|D|, \chi_D)\), by the finiteness and uniqueness theorem of Schütt ({[}Sch05, Theorem 1{]} for the finiteness ; uniqueness at our six discriminants follows from a direct LMFDB check, since each of the six relevant new spaces \(S_5^{\mathrm{new}}(|D|, \chi_D)\) contains exactly one rational eigenform). Therefore \(f_D = \theta_{\psi_D}\) and the eigenvalue identity holds.

The trace expression \(s^4 - 4 s^2 p + 2 p^2\) is the Newton power-sum \(p_4\) at \((s, n) = (s, p)\), by §2.4. \(\blacksquare\)

\hypertarget{remarks-on-the-proof}{%
\subsubsection{3.3 Remarks on the proof}\label{remarks-on-the-proof}}

\begin{enumerate}
\def\labelenumi{(\alph{enumi})}
\item
  The proof uses \emph{only} the classical theta-series identity together with the \(h_K = 1\) hypothesis (which makes every ideal principal and lets us pick a generator \(\pi\)). The rationality of the eigenvalue at split primes is automatic from \(\pi^4 + \bar\pi^4 \in \mathbb{Q}\) via the Galois action.
\item
  The \emph{uniqueness} of the rational CM newform at each of the six levels (which we use to identify \(f_D\) with \(\theta_{\psi_D}\)) is the only step that depends on the LMFDB data ; one could replace it with a direct dimension computation of \(S_5^{\mathrm{new}}(|D|, \chi_D)\) followed by a check that the rational subspace is one-dimensional (a finite computation in PARI). We sketch this in §4.4.
\item
  The choice of sign of \(a\) (i.e., between \(\pi\) and \(-\pi\), respectively between \(\pi\) and \(\bar\pi\)) is immaterial because the formula \(a^4 - 4 a^2 p + 2 p^2\) is even in \(a\) and symmetric in \(\pi, \bar\pi\).
\item
  The same argument gives, more generally, \(a_p(\theta_{\psi}) = \pi^{k-1} + \bar\pi^{k-1}\) for any class-number-\(1\) imaginary quadratic field \(K\) and any Hecke Grössencharakter \(\psi\) of infinity-type \((k - 1, 0)\) with trivial conductor. The case \(k = 5\) gives \(\pi^4 + \bar\pi^4\) ; the case \(k = 3\) gives \(\pi^2 + \bar\pi^2 = s^2 - 2p\) ; the case \(k = 7\) gives \(\pi^6 + \bar\pi^6 = s^6 - 6 s^4 p + 9 s^2 p^2 - 2 p^3\) ; etc. (See §6.1 for the generalisation to \(k \geq 7\).)
\end{enumerate}

\begin{center}\rule{0.5\linewidth}{0.5pt}\end{center}

\hypertarget{numerical-verification}{%
\subsection{4. Numerical verification}\label{numerical-verification}}

\hypertarget{the-48-pair-table}{%
\subsubsection{4.1 The 48-pair table}\label{the-48-pair-table}}

We verified Theorem A on the following \(6 \times 8\) table of \((D, p)\) pairs. The choice of eight split primes per discriminant follows the smallest split primes \(>\) ramified, with one or two skipped to keep the table balanced :

\begin{longtable}[]{@{}ll@{}}
\toprule\noalign{}
\(D\) & Split primes used \\
\midrule\noalign{}
\endhead
\bottomrule\noalign{}
\endlastfoot
\(-7\) & \(11, 23, 29, 37, 43, 53, 67, 71\) \\
\(-11\) & \(3, 5, 23, 31, 37, 47, 53, 59\) \\
\(-19\) & \(5, 7, 11, 17, 23, 43, 47, 61\) \\
\(-43\) & \(11, 13, 17, 23, 41, 47, 53, 79\) \\
\(-67\) & \(23, 29, 37, 47, 59, 71, 73, 83\) \\
\(-163\) & \(41, 43, 47, 53, 61, 79, 83, 89\) \\
\end{longtable}

All 48 pairs have \(\bigl(\tfrac{D}{p}\bigr) = +1\) (verified by PARI \texttt{kronecker(D,\ p)}).

\hypertarget{sample-table-d--67-eight-split-primes}{%
\subsubsection{\texorpdfstring{4.2 Sample table : \(D = -67\), eight split primes}{4.2 Sample table : D = -67, eight split primes}}\label{sample-table-d--67-eight-split-primes}}

For each split prime \(p\) we report : the trace \(s = a\) (from \(4p = a^2 + 67 b^2\)), the norm \(n = p\), the Newton-predicted eigenvalue \(a_p^{\mathrm{Newton}} = a^4 - 4 a^2 p + 2 p^2\), and the PARI-computed eigenvalue \(a_p^{\mathrm{PARI}}\) from \texttt{mfcoef(EB{[}1{]},\ p)} where \texttt{EB\ =\ mfeigenbasis(mfinit({[}67,\ 5,\ -67{]},\ 0))} and \texttt{EB{[}1{]}} is the unique CM eigenform (identified by \(a_q = 0\) at any inert prime, e.g.~\(q = 11\)).

\begin{longtable}[]{@{}
  >{\raggedright\arraybackslash}p{(\columnwidth - 10\tabcolsep) * \real{0.0463}}
  >{\raggedright\arraybackslash}p{(\columnwidth - 10\tabcolsep) * \real{0.3056}}
  >{\raggedright\arraybackslash}p{(\columnwidth - 10\tabcolsep) * \real{0.0833}}
  >{\raggedright\arraybackslash}p{(\columnwidth - 10\tabcolsep) * \real{0.2778}}
  >{\raggedright\arraybackslash}p{(\columnwidth - 10\tabcolsep) * \real{0.2222}}
  >{\raggedright\arraybackslash}p{(\columnwidth - 10\tabcolsep) * \real{0.0648}}@{}}
\toprule\noalign{}
\begin{minipage}[b]{\linewidth}\raggedright
\(p\)
\end{minipage} & \begin{minipage}[b]{\linewidth}\raggedright
\((a, b)\) in \(4p = a^2 + 67 b^2\)
\end{minipage} & \begin{minipage}[b]{\linewidth}\raggedright
\(s = a\)
\end{minipage} & \begin{minipage}[b]{\linewidth}\raggedright
\(\pi^4 + \bar\pi^4\) (Newton)
\end{minipage} & \begin{minipage}[b]{\linewidth}\raggedright
\(a_p(f_{-67})\) (PARI)
\end{minipage} & \begin{minipage}[b]{\linewidth}\raggedright
match
\end{minipage} \\
\midrule\noalign{}
\endhead
\bottomrule\noalign{}
\endlastfoot
23 & \((\pm 5, \pm 1)\) & \(5\) & \(625 - 2300 + 1058 = -617\) & \(-617\) & \\
29 & \((\pm 7, \pm 1)\) & \(7\) & \(2401 - 5684 + 1682 = -1601\) & \(-1601\) & \\
37 & \((\pm 9, \pm 1)\) & \(9\) & \(6561 - 11988 + 2738 = -2689\) & \(-2689\) & \\
47 & \((\pm 11, \pm 1)\) & \(11\) & \(14641 - 22748 + 4418 = -3689\) & \(-3689\) & \\
59 & \((\pm 13, \pm 1)\) & \(13\) & \(28561 - 39884 + 6962 = -4361\) & \(-4361\) & \\
71 & \((\pm 4, \pm 2)\) & \(4\) & \(256 - 4544 + 10082 = +5794\) & \(+5794\) & \\
73 & \((\pm 15, \pm 1)\) & \(15\) & \(50625 - 65700 + 10658 = -4417\) & \(-4417\) & \\
83 & \((\pm 8, \pm 2)\) & \(8\) & \(4096 - 21248 + 13778 = -3374\) & \(-3374\) & \\
\end{longtable}

All 8 match. The PARI values were computed using the script

\begin{Shaded}
\begin{Highlighting}[]
\NormalTok{default(parisize, 16*10\^{}9);}
\NormalTok{default(realprecision, 50);}
\NormalTok{G = mfinit([67, 5, {-}67], 0);}
\NormalTok{EB = mfeigenbasis(G);}
\NormalTok{\textbackslash{}\textbackslash{} identify CM form via vanishing at inert prime 11 :}
\NormalTok{for(idx=1, \#EB, print("EB", idx, "\_a\_11 = ", mfcoef(EB[idx], 11)));}
\NormalTok{\textbackslash{}\textbackslash{} The eigenform with a\_11 = 0 is the CM form ; report a\_p at split primes}
\NormalTok{for(p in [23, 29, 37, 47, 59, 71, 73, 83],}
\NormalTok{    print("a\_", p, " = ", mfcoef(EB[CM\_idx], p)));}
\end{Highlighting}
\end{Shaded}

(Full script in supplementary material \texttt{/root/crossed-cosmos/scripts/vast\_2026\_05\_10/schutt\_MULTI\_D\_optimized.py} ; PARI helper \texttt{/tmp/schutt\_MD\_D67\_w5\_p23.gp} and analogues.)

\hypertarget{verification-at-the-other-five-discriminants}{%
\subsubsection{4.3 Verification at the other five discriminants}\label{verification-at-the-other-five-discriminants}}

The same procedure was applied at \(D \in \{-7, -11, -19, -43, -163\}\). We give one entry per discriminant as a representative ; the full \(6 \times 8 = 48\)-entry table is in supplementary material.

\begin{longtable}[]{@{}
  >{\raggedright\arraybackslash}p{(\columnwidth - 12\tabcolsep) * \real{0.0952}}
  >{\raggedright\arraybackslash}p{(\columnwidth - 12\tabcolsep) * \real{0.0595}}
  >{\raggedright\arraybackslash}p{(\columnwidth - 12\tabcolsep) * \real{0.1190}}
  >{\raggedright\arraybackslash}p{(\columnwidth - 12\tabcolsep) * \real{0.0595}}
  >{\raggedright\arraybackslash}p{(\columnwidth - 12\tabcolsep) * \real{0.3571}}
  >{\raggedright\arraybackslash}p{(\columnwidth - 12\tabcolsep) * \real{0.2262}}
  >{\raggedright\arraybackslash}p{(\columnwidth - 12\tabcolsep) * \real{0.0833}}@{}}
\toprule\noalign{}
\begin{minipage}[b]{\linewidth}\raggedright
\(D\)
\end{minipage} & \begin{minipage}[b]{\linewidth}\raggedright
\(p\)
\end{minipage} & \begin{minipage}[b]{\linewidth}\raggedright
\((a, b)\)
\end{minipage} & \begin{minipage}[b]{\linewidth}\raggedright
\(s\)
\end{minipage} & \begin{minipage}[b]{\linewidth}\raggedright
Newton \(\pi^4 + \bar\pi^4\)
\end{minipage} & \begin{minipage}[b]{\linewidth}\raggedright
PARI \(a_p(f_D)\)
\end{minipage} & \begin{minipage}[b]{\linewidth}\raggedright
match
\end{minipage} \\
\midrule\noalign{}
\endhead
\bottomrule\noalign{}
\endlastfoot
\(-7\) & 11 & \((\pm 4, \pm 2)\) & \(4\) & \(256 - 704 + 242 = -206\) & \(-206\) & \\
\(-11\) & 3 & \((\pm 1, \pm 1)\) & \(1\) & \(1 - 12 + 18 = +7\) & \(+7\) & \\
\(-19\) & 5 & \((\pm 1, \pm 1)\) & \(1\) & \(1 - 20 + 50 = +31\) & \(+31\) & \\
\(-43\) & 11 & \((\pm 1, \pm 1)\) & \(1\) & \(1 - 44 + 242 = +199\) & \(+199\) & \\
\(-67\) & 23 & \((\pm 5, \pm 1)\) & \(5\) & \(625 - 2300 + 1058 = -617\) & \(-617\) & \\
\(-163\) & 41 & \((\pm 1, \pm 1)\) & \(1\) & \(1 - 164 + 3362 = +3199\) & \(+3199\) & \\
\end{longtable}

(For \(D = -7\), \(p = 11\) : \(44 = a^2 + 7 b^2\) has solution \((4, 2)\) : \(16 + 28 = 44\) (the candidate \((3, 1)\) gives \(9 + 7 = 16 \ne 44\), so is invalid). For \(D = -11\), \(p = 3\) : \(12 = a^2 + 11 b^2\) has \((1, 1)\) : \(1 + 11 = 12\) . For \(D = -19\), \(p = 5\) : \(20 = a^2 + 19 b^2\) has \((1, 1)\) : \(1 + 19 = 20\) . For \(D = -43\), \(p = 11\) : \(44 = a^2 + 43 b^2\) has \((1, 1)\) : \(1 + 43 = 44\) . For \(D = -67\), \(p = 23\) : \(92 = a^2 + 67 b^2\) has \((5, 1)\) : \(25 + 67 = 92\) . For \(D = -163\), \(p = 41\) : \(164 = a^2 + 163 b^2\) has \((1, 1)\) : \(1 + 163 = 164\) .)

\hypertarget{identification-of-the-cm-eigenform-within-s_5mathrmnewd-chi_d}{%
\subsubsection{\texorpdfstring{4.4 Identification of the CM eigenform within \(S_5^{\mathrm{new}}(|D|, \chi_D)\)}{4.4 Identification of the CM eigenform within S\_5\^{}\{\textbackslash mathrm\{new\}\}(\textbar D\textbar, \textbackslash chi\_D)}}\label{identification-of-the-cm-eigenform-within-s_5mathrmnewd-chi_d}}

A subtle point : at \(D = -67\) the new space \(S_5^{\mathrm{new}}(67, \chi_{-67})\) has dimension \(> 1\), with two Galois orbits of newforms. The unique rational CM eigenform is identified as \texttt{EB{[}i{]}} with \(a_q = 0\) at any inert prime \(q\) (CM signature). At inert prime \(q = 11\), only one eigenform vanishes ; this is the CM form. The other eigenform (the ``non-CM rational'' or ``non-rational orbit'', depending on the discriminant) has nonzero \(a_{11}\) and is \emph{not} relevant for Theorem A.

Dimensions of the new spaces and the CM rational eigenforms (LMFDB cross-check) :

\begin{longtable}[]{@{}
  >{\raggedright\arraybackslash}p{(\columnwidth - 6\tabcolsep) * \real{0.0816}}
  >{\raggedright\arraybackslash}p{(\columnwidth - 6\tabcolsep) * \real{0.4286}}
  >{\raggedright\arraybackslash}p{(\columnwidth - 6\tabcolsep) * \real{0.3265}}
  >{\raggedright\arraybackslash}p{(\columnwidth - 6\tabcolsep) * \real{0.1633}}@{}}
\toprule\noalign{}
\begin{minipage}[b]{\linewidth}\raggedright
\(D\)
\end{minipage} & \begin{minipage}[b]{\linewidth}\raggedright
\(\dim S_5^{\mathrm{new}}(|D|, \chi_D)\)
\end{minipage} & \begin{minipage}[b]{\linewidth}\raggedright
Number of rational eigenforms
\end{minipage} & \begin{minipage}[b]{\linewidth}\raggedright
CM form label
\end{minipage} \\
\midrule\noalign{}
\endhead
\bottomrule\noalign{}
\endlastfoot
\(-7\) & \(1\) & \(1\) & \(7.5.b.a\) \\
\(-11\) & \(1\) & \(1\) & \(11.5.b.a\) \\
\(-19\) & \(1\) & \(1\) & \(19.5.b.a\) \\
\(-43\) & \(1\) & \(1\) & \(43.5.b.a\) \\
\(-67\) & \(2\) & \(1\) & \(67.5.b.a\) \\
\(-163\) & (large) & \(1\) & \(163.5.b.a\) \\
\end{longtable}

The dimension of \(S_5^{\mathrm{new}}(67, \chi_{-67})\) as computed by PARI \texttt{mfdim(mfinit({[}67,\ 5,\ -67{]},\ 0))} is in fact (for the Atkin--Lehner-character convention of LMFDB) \(2\), with the two Galois orbits being the rational CM form \(67.5.b.a\) (dimension \(1\)) and a single Galois orbit of dimension \(1\) (giving total dim \(2\)). The fact that there are \emph{two} rational eigenforms --- only one of them CM --- at \(D = -67\) requires the CM identification step (vanishing at inert primes). At \(D = -7, -11, -19, -43, -163\) the situation is simpler : the new space has dimension \(1\) and the unique eigenform is automatically CM.

\hypertarget{computational-notes}{%
\subsubsection{4.5 Computational notes}\label{computational-notes}}

The script \texttt{/root/crossed-cosmos/scripts/vast\_2026\_05\_10/schutt\_MULTI\_D\_optimized.py} (8 D × 8 primes = 64 tasks, 32-parallel on a 96-core EPYC machine, \(\sim 5\)/15min wall-time) computes all relevant \((D, p)\) pairs and verifies the Newton match. The overall outcome is \texttt{64/64\ OK}; the relevant subset to Theorem A is the \(48\) pairs with \(D \in \{-7, -11, -19, -43, -67, -163\}\), all of which give Newton \(=\) PARI to 50-digit precision.

For the purpose of this paper we restrict the table to the six \(h_K = 1\) discriminants ; the script's two extra discriminants (\(D = -84\) with \(h_K = 4\) and \(D = -148\) with \(h_K = 2\)) gave Newton-non-match exactly as expected, since their CM eigenvalues live in a degree-\(h_K\) extension (cf.~§6.2).

\begin{center}\rule{0.5\linewidth}{0.5pt}\end{center}

\hypertarget{the-hodge-theoretic-interpretation-sketch}{%
\subsection{5. The Hodge-theoretic interpretation : sketch}\label{the-hodge-theoretic-interpretation-sketch}}

This section sketches the conjectural interpretation of Theorem A in terms of the Hodge conjecture for the absolute eight-fold self-product \((E_K)^8 / \mathbb{Q}\) of the canonical CM elliptic curve \(E_K\) over \(\mathbb{Q}\) at \(D \in \{-7, -11, -19, -43, -67, -163\}\). The exposition follows the lines of {[}Sch08{]} (for the Picard-rank-\(20\) K3 setting) and the more recent Costa--Elsenhans--Jahnel--Voight 2025 framework {[}CEJV25{]} (for the Kuga--Satake correspondence in the CM K3 setting), but adapted to our weight-\(5\) setting where the structurally correct host is the eight-fold self-product of the elliptic curve, not the K3 surface itself.

The discussion in this section is \emph{informal and conjectural} ; the rigorous content of the paper is Theorem A, which is established in §3 and verified in §4. Section 5 is intended to motivate the eight-fold framework and to identify the explicit \(5\)-dimensional Hecke eigencomponent in \(H^8((E_K)^8, \mathbb{Q})\) that \emph{would} host the Newton-identity eigenvalues if the Hodge conjecture for this \(8\)-fold were known.

\hypertarget{the-canonical-cm-elliptic-curve-e_k-over-mathbbq}{%
\subsubsection{\texorpdfstring{5.1 The canonical CM elliptic curve \(E_K\) over \(\mathbb{Q}\)}{5.1 The canonical CM elliptic curve E\_K over \textbackslash mathbb\{Q\}}}\label{the-canonical-cm-elliptic-curve-e_k-over-mathbbq}}

For each \(D \in \{-7, -11, -19, -43, -67, -163\}\), the class number \(h_K = 1\) implies that the CM elliptic curve \(E_K\) with \(\mathrm{End}_{\bar{\mathbb{Q}}}(E_K) = \mathcal{O}_K\) is defined over \(\mathbb{Q}\) (up to twist), with the unique \(j\)-invariant \(j(E_K) \in \mathbb{Z}\) being the Heegner singular modulus. Explicit values :
- \(D = -7\) : \(j(E_K) = -3375 = -3^3 \cdot 5^3\)
- \(D = -11\) : \(j(E_K) = -32768 = -2^{15}\)
- \(D = -19\) : \(j(E_K) = -884736 = -2^{15} \cdot 3^3\)
- \(D = -43\) : \(j(E_K) = -884736000 = -2^{18} \cdot 3^3 \cdot 5^3\)
- \(D = -67\) : \(j(E_K) = -147197952000 = -2^{15} \cdot 3^3 \cdot 5^3 \cdot 11^3\)
- \(D = -163\) : \(j(E_K) = -262537412640768000 = -2^{18} \cdot 3^3 \cdot 5^3 \cdot 23^3 \cdot 29^3\)

The conductor of \(E_K\) is \(|D|^2\) (the unique bad reduction is at \(|D|\), and is of additive type).

\hypertarget{the-galois-representation-on-h1e_k-mathbbq_ell}{%
\subsubsection{\texorpdfstring{5.2 The Galois representation on \(H^1(E_K, \mathbb{Q}_\ell)\)}{5.2 The Galois representation on H\^{}1(E\_K, \textbackslash mathbb\{Q\}\_\textbackslash ell)}}\label{the-galois-representation-on-h1e_k-mathbbq_ell}}

For \(\ell \neq p\) the Tate module \(T_\ell E_K = H^1(E_K, \mathbb{Q}_\ell)^\vee\) is a \(2\)-dimensional \(\mathbb{Q}_\ell\)-vector space carrying a continuous action of the absolute Galois group \(G_\mathbb{Q}\). Restricted to \(G_K\), this action splits as \(\psi_E \oplus \bar\psi_E\) where \(\psi_E\) is the canonical Hecke Grössencharakter of \(K\) of infinity-type \((1, 0)\). Equivalently, \(H^1(E_K, \mathbb{Q}_\ell) = \mathrm{Ind}_{G_K}^{G_\mathbb{Q}}(\psi_E) \otimes \mathbb{Q}_\ell\).

At a prime \(p\) split in \(K\), \(p\,\mathcal{O}_K = \mathfrak{p}\,\bar{\mathfrak{p}}\), the Frobenius \(\mathrm{Frob}_p\) acts on \(H^1(E_K, \mathbb{Q}_\ell)\) with eigenvalues \(\psi_E(\mathfrak{p}) = \pi\) and \(\psi_E(\bar{\mathfrak{p}}) = \bar\pi\) (with \(\pi\bar\pi = p\) by Weil). At an inert prime, \(\mathrm{Frob}_p\) acts with eigenvalues \(\pm i\sqrt p\) (trace zero, the ``CM signature'').

\hypertarget{the-four-fold-absolute-self-product-y_4-e_k4-mathbbq-and-its-h4}{%
\subsubsection{\texorpdfstring{5.3 The four-fold absolute self-product \(Y_4 = (E_K)^4 / \mathbb{Q}\) and its \(H^4\)}{5.3 The four-fold absolute self-product Y\_4 = (E\_K)\^{}4 / \textbackslash mathbb\{Q\} and its H\^{}4}}\label{the-four-fold-absolute-self-product-y_4-e_k4-mathbbq-and-its-h4}}

Define \(Y_4 := E_K \times_\mathbb{Q} E_K \times_\mathbb{Q} E_K \times_\mathbb{Q} E_K\) (four copies) over \(\mathrm{Spec}\,\mathbb{Q}\). By Künneth,
\[
H^4(Y_4, \mathbb{Q}_\ell) \;=\; \wedge^4 \bigl(H^1(E_K, \mathbb{Q}_\ell)^{\oplus 4}\bigr) \;\cong\; \mathbb{Q}_\ell^{\binom{8}{4}} \;=\; \mathbb{Q}_\ell^{70},
\]
of weight \(4\). The Hodge filtration on \(H^4(Y_4, \mathbb{C})\) has Hodge numbers \(h^{p, q}(Y_4) = \binom{4}{p} \binom{4}{q}\) (for \(p + q = 4\)), summing to \(70\).

At a split prime \(p\,\mathcal{O}_K = \mathfrak{p}\,\bar{\mathfrak{p}}\), the Frobenius \(\mathrm{Frob}_p\) acts on \(H^4(Y_4, \mathbb{Q}_\ell)\) with eigenvalues
\[
\bigl\{ \pi^k\,\bar\pi^{4 - k} : 0 \leq k \leq 4 \bigr\}
\]
each with multiplicity \(\binom{4}{k}^2\) (combinatorial : choose \(k\) of \(4\) tensor factors to contribute \(\pi\) from the ``pi-half'' of \(H^1\), and similarly \(4 - k\) for \(\bar\pi\)).

The \(5\)-tuple of eigenvalues \(\{\pi^4, \pi^3 \bar\pi, \pi^2 \bar\pi^2, \pi \bar\pi^3, \bar\pi^4\}\) is exactly that of \(\mathrm{Sym}^4 \rho_{\psi_E}\) (cf.~§5.4(b)) ; the \(S_4\)-invariant subspace of \(H^4(Y_4)\) realises this \(5\)-dim Galois representation as a sub-Hecke-eigencomponent.

(\emph{Historical aside.} An alternative --- and ultimately mistaken --- framing in earlier drafts of this paper considered the \emph{eight-fold} self-product \((E_K)^8\) of dimension \(\binom{16}{8} = 12\,870\) for the embedding of \(\rho_{f_D}\). The mistake was a Tate-twist accounting error that confused the weight of \(\rho_{f_D}\) --- which is \(4\) --- with the weight of \(H^8((E_K)^8)\) --- which is \(8\). The cleanest and most natural host is the four-fold \(Y_4 = (E_K)^4\) with \(H^4\) of weight \(4\), requiring no Tate twist.)

\hypertarget{two-galois-reps-to-embed-rho_f_d-vs-mathrmsym4-rho_psi_e}{%
\subsubsection{\texorpdfstring{5.4 Two Galois reps to embed : \(\rho_{f_D}\) vs \(\mathrm{Sym}^4 \rho_{\psi_E}\)}{5.4 Two Galois reps to embed : \textbackslash rho\_\{f\_D\} vs \textbackslash mathrm\{Sym\}\^{}4 \textbackslash rho\_\{\textbackslash psi\_E\}}}\label{two-galois-reps-to-embed-rho_f_d-vs-mathrmsym4-rho_psi_e}}

There are two distinct \(2\)-dim and \(5\)-dim Galois representations to consider, naturally associated to \(f_D\) and the underlying CM elliptic curve \(E_K\). We disambiguate them carefully.

\begin{enumerate}
\def\labelenumi{(\alph{enumi})}
\item
  \textbf{The \(2\)-dim Galois representation \(\rho_{f_D}\) attached to \(f_D\)}, of motivic weight \(k - 1 = 4\) (since \(f_D\) is weight \(k = 5\)). Its Frobenius eigenvalues at a split prime \(p\) are \(\alpha = \psi_D(\mathfrak{p}) = \pi^4\) and \(\beta = \psi_D(\bar{\mathfrak{p}}) = \bar\pi^4\), with \(\alpha\beta = (\pi\bar\pi)^4 = p^4\) matching the expected \(\det \rho_{f_D}(\mathrm{Frob}_p) = \chi_D(p) p^{k - 1} = p^4\). The trace \(\alpha + \beta = \pi^4 + \bar\pi^4 = a_p(f_D)\) is the eigenvalue computed in Theorem A.
\item
  \textbf{The \(5\)-dim Galois representation \(\mathrm{Sym}^4 \rho_{\psi_E}\)}, where \(\rho_{\psi_E}\) is the \(2\)-dim rep on \(H^1(E_K, \mathbb{Q}_\ell)^\vee\) attached to the \emph{weight-\(1\)} Hecke Grössencharakter \(\psi_E\) of infinity-type \((1, 0)\) (so \(\rho_{\psi_E}\) has motivic weight \(1\)). Its Frobenius eigenvalues at split \(p\) are \(\alpha_E = \pi\), \(\beta_E = \bar\pi\) (with \(\alpha_E\beta_E = p\)). Then \(\mathrm{Sym}^4 \rho_{\psi_E}\) has \(5\) eigenvalues \(\{\alpha_E^4, \alpha_E^3\beta_E, \alpha_E^2\beta_E^2, \alpha_E\beta_E^3, \beta_E^4\} = \{\pi^4, \pi^3\bar\pi, \pi^2\bar\pi^2, \pi\bar\pi^3, \bar\pi^4\}\), of total motivic weight \(4 \cdot 1 = 4\) (each eigenvalue has absolute value \(p^2\)).
\end{enumerate}

\textbf{Key relation.} The trace of \(\mathrm{Sym}^4 \rho_{\psi_E}\) at split \(p\) is
\[
\mathrm{Tr}\,\mathrm{Sym}^4 \rho_{\psi_E}(\mathrm{Frob}_p) \;=\; \pi^4 + \pi^3\bar\pi + \pi^2\bar\pi^2 + \pi\bar\pi^3 + \bar\pi^4 \;=\; \frac{\pi^5 - \bar\pi^5}{\pi - \bar\pi} \cdot (\text{unit factor}).
\]
This is \emph{not} the same as \(a_p(f_D) = \pi^4 + \bar\pi^4\) ! The two are related but distinct.

The relation between (a) and (b) is :
\[
\rho_{f_D} \;=\; \rho_{\psi_E^4} \;=\; \rho_{\psi_E}^{\otimes 4}\bigr|_{\text{symmetric}} \;-\; (\text{lower Sym pieces}).
\]
More precisely, the \(4\)-th symmetric power \(\mathrm{Sym}^4 \rho_{\psi_E}\) decomposes over \(K\) as a sum of \(5\) characters, of which the two ``extreme'' characters \(\psi_E^4\) and \(\bar\psi_E^4\) assemble to the Galois rep \(\rho_{f_D} = \mathrm{Ind}_{G_K}^{G_\mathbb{Q}}(\psi_E^4)\) on the \(2\)-dimensional weight-\(4\) subspace. The other three characters \(\{\psi_E^3 \bar\psi_E, \psi_E^2 \bar\psi_E^2, \psi_E \bar\psi_E^3\}\) assemble to a \(3\)-dimensional sub-rep of \(\mathrm{Sym}^4 \rho_{\psi_E}\) that does \emph{not} correspond to \(f_D\).

\hypertarget{embedding-rho_f_d-into-h4e_k4-2}{%
\subsubsection{\texorpdfstring{5.5 Embedding \(\rho_{f_D}\) into \(H^4((E_K)^4)(-2)\)}{5.5 Embedding \textbackslash rho\_\{f\_D\} into H\^{}4((E\_K)\^{}4)(-2)}}\label{embedding-rho_f_d-into-h4e_k4-2}}

The \(2\)-dim weight-\(4\) rep \(\rho_{f_D}\) can be naturally embedded inside the cohomology of the \(4\)-fold self-product \(Y_4 := (E_K)^4 / \mathbb{Q}\). By Künneth,
\[
H^4(Y_4, \mathbb{Q}_\ell) \;=\; \wedge^4 \bigl(H^1(E_K, \mathbb{Q}_\ell)^{\oplus 4}\bigr),
\]
of dimension \(\binom{8}{4} = 70\), weight \(4\). Its \(S_4\)-invariant part contains \(\mathrm{Sym}^4 H^1(E_K, \mathbb{Q}_\ell)\) as a \(5\)-dimensional subspace, which (by §5.4(b)) hosts the \(5\)-character decomposition of \(\mathrm{Sym}^4 \rho_{\psi_E}\). The two ``extreme'' characters \(\psi_E^4 \oplus \bar\psi_E^4 = \rho_{f_D}\) form a \(2\)-dimensional Galois subrepresentation. \emph{No Tate twist is needed for this embedding} --- both sides are weight \(4\).

This is the \emph{cleanest} Hodge-theoretic interpretation of Theorem A : the eigenvalues \(a_p(f_D) = \pi^4 + \bar\pi^4\) of Theorem A appear as the trace of Frobenius on a \(2\)-dim subspace of \(\mathrm{Sym}^4 H^1((E_K)^4, \mathbb{Q}_\ell)\) corresponding to \(\rho_{f_D}\).

\hypertarget{the-full-hodge-class-question}{%
\subsubsection{5.6 The full Hodge-class question}\label{the-full-hodge-class-question}}

The Hodge-theoretic open question is whether the \(2\)-dim subspace \(\rho_{f_D} \subset H^4(Y_4, \mathbb{Q}_\ell)\) is the étale realisation of an \emph{algebraic Hodge class} on \(Y_4 = (E_K)^4 / \mathbb{Q}\) : explicitly, whether the Hodge realisation of \(\rho_{f_D}\) in \(H^{2, 2}(Y_4, \mathbb{C}) \cap H^4(Y_4, \mathbb{Q})\) is the cohomology class of an algebraic cycle of dimension \(2\) on \(Y_4\).

For self-products of CM elliptic curves of class number \(1\), this Hodge-class algebraicity \emph{can} be derived from the rich endomorphism algebra \(\mathcal{O}_K \otimes \mathcal{O}_K \otimes \mathcal{O}_K \otimes \mathcal{O}_K\) acting on \(H^4(Y_4)\), by combining Tankeev's theorem (for \(H^2\) of CM abelian varieties) with the standard Künneth-formula construction. We do not give the explicit construction here ; see {[}Tot19, §3{]} for the relevant techniques and {[}Sch05, §3{]} for the analogous statement at lower-dimensional CM K3 surfaces.

For \emph{higher} symmetric powers (such as the full \(5\)-dim \(\mathrm{Sym}^4 \rho_{\psi_E}\), which goes beyond \(\rho_{f_D}\)), the relevant host is \(H^4((E_K)^4)\) as well (since both have weight \(4\)), but the algebraicity of the full \(5\)-dim Hecke eigencomponent is more subtle and relates to the deeper aspects of the algebraicity programme.

\hypertarget{conjectural-hodge-class-statement}{%
\subsubsection{5.7 Conjectural Hodge-class statement}\label{conjectural-hodge-class-statement}}

\begin{quote}
\textbf{Conjecture 5.7} (Hodge-class algebraicity for \(\rho_{f_D}\) on \((E_K)^4\)).
Let \(D \in \{-7, -11, -19, -43, -67, -163\}\) and let \(E_K\) be the canonical CM elliptic curve with \(\mathrm{End}_{\bar{\mathbb{Q}}}(E_K) = \mathcal{O}_{K_D}\). Let \(Y_4 := (E_K)^4 / \mathbb{Q}\). The \(2\)-dim Hecke eigencomponent of \(H^4(Y_4, \mathbb{Q})\) corresponding to the rep \(\rho_{f_D} = \mathrm{Ind}_{G_K}^{G_\mathbb{Q}}(\psi_E^4)\) is the étale realisation of an algebraic Hodge class : there exists an explicit algebraic cycle \(Z_D \subset Y_4\) of dimension \(2\) whose cohomology class generates this \(2\)-dim subspace.
\end{quote}

This conjecture is \textbf{expected to follow} from Tankeev's theorem on Hodge cycles of CM abelian varieties (see {[}Tot19, §3.2{]}) combined with the explicit endomorphism action \(\mathcal{O}_{K_D}^{\otimes 4}\) on \(H^4(Y_4)\), but a \emph{detailed construction} of \(Z_D\) for each \(D\) is not in the literature to the author's knowledge. Theorem A provides the \emph{eigenvalue evidence} (the trace of Frobenius on the conjectural algebraic-cycle \(Z_D\) at split primes is exactly \(a_p(f_D) = \pi^4 + \bar\pi^4\)).

The (much stronger) Hodge conjecture for \(H^4\) of all \(4\)-folds --- which would settle Conjecture 5.7 unconditionally --- is open. For the \emph{specific} CM context, however, Conjecture 5.7 is \emph{much weaker} than the full Hodge conjecture and is widely expected to hold. We do not prove it here.

\hypertarget{relation-to-schuxfctts-k3-framework-and-the-costaejv-kugasatake-correspondence}{%
\subsubsection{5.8 Relation to Schütt's K3 framework and the Costa--EJV Kuga--Satake correspondence}\label{relation-to-schuxfctts-k3-framework-and-the-costaejv-kugasatake-correspondence}}

The relation to existing literature on CM K3 surfaces and their eight-fold self-products is as follows :
- \textbf{Schütt's K3 framework} ({[}Sch05{]}, {[}Sch08{]}) classifies CM K3 surfaces of Picard rank \(20\) over \(\mathbb{Q}\) and their \(H^2\)-modularity in terms of weight-\(3\) CM newforms. The transcendental lattice \(T(X)\) of such a K3 has rank \(2\) and motivic weight \(2\), hosting a \(2\)-dimensional Galois representation isomorphic to a weight-\(3\) CM newform's Galois representation. The connection to \emph{higher} symmetric powers (\(\mathrm{Sym}^n\) for \(n \geq 2\)) is \emph{not} via \(H^2(X)\) directly but via higher self-products of \(X\) and twists.
- \textbf{Costa--Elsenhans--Jahnel--Voight 2025} ({[}CEJV25{]}) extends Schütt's framework to CM K3 surfaces over higher base fields, expressing the transcendental motive \(T(X)\) as a \emph{wedge subsummand} of \(\wedge^2 H^1(A)\) for an auxiliary CM abelian threefold \(A\) via the \emph{Kuga--Satake} correspondence. This is a different construction from the Kuga--Sato modular variety.
- \textbf{Our framework} uses neither K3 surfaces nor Kuga--Satake correspondence : we directly take the four-fold self-product \((E_K)^4\) of the CM elliptic curve \(E_K\) over \(\mathbb{Q}\) and identify the \(2\)-dimensional Hecke eigencomponent of its middle cohomology \(H^4((E_K)^4, \mathbb{Q})\) corresponding to \(\rho_{f_D}\). This is the cleanest framework for stating Theorem A (which is a pure eigenvalue identity, not a Hodge claim) and for stating its conjectural Hodge-class refinement (Conjecture 5.7).

\begin{center}\rule{0.5\linewidth}{0.5pt}\end{center}

\hypertarget{generalisations-and-open-questions}{%
\subsection{6. Generalisations and open questions}\label{generalisations-and-open-questions}}

\hypertarget{higher-weight-k-geq-7}{%
\subsubsection{\texorpdfstring{6.1 Higher weight \(k \geq 7\)}{6.1 Higher weight k \textbackslash geq 7}}\label{higher-weight-k-geq-7}}

The same theta-series argument generalises immediately. For each \(D \in \{-7, -11, -19, -43, -67, -163\}\) and each odd \(k \geq 3\) for which a Hecke Grössencharakter \(\psi_D^{(k - 1)}\) of infinity-type \((k - 1, 0)\) exists with trivial conductor (which holds for all \(k \geq 3\) in the \(h_K = 1\) setting), the associated theta-series \(\theta_{\psi_D^{(k-1)}}\) is a weight-\(k\) CM newform in \(S_k^{\mathrm{new}}(|D|, \chi_D)\) with rational Fourier coefficients given by
\[
a_p \;=\; \pi^{k - 1} + \bar\pi^{k - 1}
\]
at split primes \(p\,\mathcal{O}_K = \mathfrak{p}\,\bar{\mathfrak{p}}\). The Newton recursion gives the explicit polynomial in \((s, p) = (\pi + \bar\pi, p)\) at any \(k\). For example :

\begin{itemize}
\tightlist
\item
  \(k = 7\) : \(a_p = s^6 - 6 s^4 p + 9 s^2 p^2 - 2 p^3\)
\item
  \(k = 9\) : \(a_p = s^8 - 8 s^6 p + 20 s^4 p^2 - 16 s^2 p^3 + 2 p^4\)
\item
  \(k = 11\) : \(a_p = s^{10} - 10 s^8 p + 35 s^6 p^2 - 50 s^4 p^3 + 25 s^2 p^4 - 2 p^5\)
\end{itemize}

(General formula : \(p_k = \sum_{j = 0}^{\lfloor k/2 \rfloor} (-1)^j \frac{k}{k - j} \binom{k - j}{j} s^{k - 2j} p^j\), where the leading sign is positive for even \(k\) and the constant term is \((-1)^{k/2} \cdot 2 \cdot p^{k/2}\) when \(k\) is even.)

A complete tabulation of \(a_p\) at the same \(48\) pairs at weight \(7, 9, 11, 13\) would extend the present table to \(192\) entries ; we leave this as a routine exercise.

\hypertarget{the-case-d--8-i.e.-k-mathbbqsqrt-2-and-the-d-equiv-0-pmod-4-subtlety}{%
\subsubsection{\texorpdfstring{6.2 The case \(D = -8\) (i.e., \(K = \mathbb{Q}(\sqrt{-2})\)) and the \(D \equiv 0 \pmod 4\) subtlety}{6.2 The case D = -8 (i.e., K = \textbackslash mathbb\{Q\}(\textbackslash sqrt\{-2\})) and the D \textbackslash equiv 0 \textbackslash pmod 4 subtlety}}\label{the-case-d--8-i.e.-k-mathbbqsqrt-2-and-the-d-equiv-0-pmod-4-subtlety}}

The discriminant \(D = -8\) corresponds to \(K = \mathbb{Q}(\sqrt{-2})\), of class number \(1\) but \(D \equiv 0 \pmod 4\). The ring of integers is \(\mathcal{O}_K = \mathbb{Z}[\sqrt{-2}]\) (not \(\mathbb{Z}[(1 + \sqrt{-2})/2]\)), and \(\pi\) is parameterised as \(\pi = a + b\sqrt{-2}\) with \(a, b \in \mathbb{Z}\), giving \(N(\pi) = a^2 + 2 b^2\) and \(\mathrm{tr}(\pi) = 2a\). The Newton recursion thus uses \(s = 2a\), \(n = p = a^2 + 2 b^2\), giving
\[
a_p \;=\; (2a)^4 - 4 (2a)^2 p + 2 p^2 \;=\; 16 a^4 - 16 a^2 p + 2 p^2.
\]
At \(D = -8\), \(p = 3\) (smallest split prime) : \(3 = a^2 + 2 b^2\) has \((a, b) = (\pm 1, \pm 1)\). Then \(s = \pm 2\), \(a_3 = 16 - 48 + 18 = -14\). Cross-check : LMFDB form \(8.5.b.a\) has \(a_3 = -14\).

Theorem A thus extends to \(D = -8\) with the substitution \(s \to 2a\) (and the form labelled \(8.5.b.a\)). The general statement is :

\begin{quote}
\textbf{Theorem A}\('\) : For \(K = \mathbb{Q}(\sqrt D)\) with \(h_K = 1\) and \(|\mathcal{O}_K^\times| = 2\), regardless of \(D \bmod 4\), with \(\pi \in \mathcal{O}_K\) a generator of a prime above split \(p\), the eigenvalue of the canonical weight-\(5\) CM newform \(f_D\) at \(p\) is \(\pi^4 + \bar\pi^4\).
\end{quote}

The formula \(s^4 - 4 s^2 p + 2 p^2\) is then unconditional in \(\mathbb{Z}\) ; the only difference between the \(D \equiv 1\) and \(D \equiv 0 \pmod 4\) cases is the parameterisation of \(\pi\).

\hypertarget{the-case-d--3-and-d--4-mathcalo_ktimes-2}{%
\subsubsection{\texorpdfstring{6.3 The case \(D = -3\) and \(D = -4\) (\(|\mathcal{O}_K^\times| > 2\))}{6.3 The case D = -3 and D = -4 (\textbar\textbackslash mathcal\{O\}\_K\^{}\textbackslash times\textbar{} \textgreater{} 2)}}\label{the-case-d--3-and-d--4-mathcalo_ktimes-2}}

For \(D \in \{-3, -4\}\), the unit group \(\mathcal{O}_K^\times\) has order \(6\) (for \(D = -3\), with \(\mathcal{O}_K^\times = \langle \zeta_6 \rangle\)) or \(4\) (for \(D = -4\), with \(\mathcal{O}_K^\times = \langle i \rangle\)). The Hecke Grössencharakter \(\psi_K^{(4)}\) of infinity-type \((4, 0)\) then has nontrivial restriction to the units, and the theta-series identity must be modified to account for this. For \(D = -4\) : the Grössencharakter satisfies \(\psi(i \pi) = i^4 \psi(\pi) = \psi(\pi)\), so the units act trivially on the eigenvalue and Theorem A holds \emph{unchanged}. For \(D = -3\) : \(\psi(\zeta_6 \pi) = \zeta_6^4 \psi(\pi) = \zeta_3 \psi(\pi)\), so the units do \emph{not} act trivially ; the eigenvalue \(\pi^4 + \bar\pi^4\) depends on the choice of representative \(\pi\) in its \(\zeta_6\)-orbit, and there is \emph{no} canonical rational Hecke eigenvalue at general split primes.

A more careful statement at \(D = -3\) : the canonical weight-\(5\) CM newform for \(K = \mathbb{Q}(\sqrt{-3})\) corresponds to a Grössencharakter \(\psi\) satisfying \(\psi(u \pi) = u^? \psi(\pi)\) for an exponent \(?\) chosen to give rational eigenvalues. The valid exponent is \(4 + 6\,m\) for any \(m \in \mathbb{Z}\), giving infinity-type \((4 + 6m, 0)\). The minimal choice \(m = 0\) (infinity-type \((4, 0)\)) gives a non-trivial \(\zeta_6\)-action ; the minimal choice for \emph{trivial} unit action is \(m = 1\) (infinity-type \((10, 0)\)), giving a weight-\(11\) CM newform. Hence the weight-\(5\) statement does not extend cleanly to \(D = -3\).

\hypertarget{the-case-h_k-1}{%
\subsubsection{\texorpdfstring{6.4 The case \(h_K > 1\)}{6.4 The case h\_K \textgreater{} 1}}\label{the-case-h_k-1}}

For \(h_K > 1\), the \(\theta_\psi\) identity at split primes \(p\,\mathcal{O}_K = \mathfrak{p}\,\bar{\mathfrak{p}}\) becomes
\[
a_p(\theta_\psi) \;=\; \psi(\mathfrak{p}) + \psi(\bar{\mathfrak{p}}),
\]
where \(\psi(\mathfrak{p})\) is \emph{not} simply \(\pi^4\) for any \(\pi \in \mathcal{O}_K\) (since \(\mathfrak{p}\) is not principal in general). Instead, \(\psi(\mathfrak{p})\) is an algebraic number in the class field \(H_K\) of \(K\), of degree \(h_K\) over \(K\). The Galois orbit of \(\theta_\psi\) under \(\mathrm{Gal}(H_K / K)\) has size \(h_K\), and the eigenvalues at split primes lie in \(\mathbb{Q}(\zeta_{h_K} \psi(\mathfrak{p}))\), generically of degree \(h_K\) over \(\mathbb{Q}\).

For \(h_K = 2\) (e.g.~\(D = -15\)), the eigenvalues lie in a quadratic extension of \(\mathbb{Q}\) ; the analogue of Theorem A would be a \emph{quadratic-irrational} version where the right-hand side is \(\sigma_1(\pi)^4 + \sigma_2(\pi)^4\) for two embeddings \(\sigma_1, \sigma_2 : H_K \hookrightarrow \mathbb{C}\).

For \(h_K = 4\) biquadratic (e.g.~\(D = -84\)), the situation is more intricate ; the elementary 2-abelian rigidity framework of {[}Pap10{]} (equivalently, the M142-NORM-rationality framework in the larger arithmetic-physics manuscript) provides a partial generalisation in the form of a \emph{symmetric-polynomial-of-Galois-orbit} identity, but this requires a separate paper.

\hypertarget{open-questions}{%
\subsubsection{6.5 Open questions}\label{open-questions}}

\begin{enumerate}
\def\labelenumi{(\alph{enumi})}
\item
  \textbf{Algebraicity of the \(H^{2, 2}\) Hodge class on \((E_K)^4\)} : Conjecture 5.7 (the Hodge-class refinement of Theorem A) is widely expected by Tankeev's framework but lacks an explicit construction. A \emph{direct construction} of the algebraic cycle \(Z_D \subset (E_K)^4\) of dimension \(2\) representing the canonical Hecke eigencomponent corresponding to \(\rho_{f_D}\) would be a major contribution and is the natural next step after the present paper. The cycle \(Z_D\) should be expressible in terms of the explicit endomorphism action of \(\mathcal{O}_{K_D}\) on each factor of \((E_K)^4\), but writing it down combinatorially requires care.
\item
  \textbf{Generalisation to \(h_K = 4\) biquadratic} : the elementary \(2\)-abelian rigidity framework of {[}Pap10{]} suggests a symmetric-polynomial generalisation of Theorem A at \(h_K = 4\). Explicit statement is open.
\item
  \textbf{Connection to L-values} : Theorem A gives a closed-form expression for \(a_p\) at split primes, which via the Euler product gives the local L-factor at \(p\). The associated central L-value \(L(f_D, 5/2)\) has a Damerell-type closed form involving periods of \(E_K\) ({[}Yag82{]}, {[}Sch10{]}). The relation between Theorem A's eigenvalue identity and the closed-form L-value is the content of the M142 hierarchy of \(10\) proven theorems in the larger arithmetic-physics manuscript ; the abstract relationship is via the symmetric-square L-function \(L(\mathrm{Sym}^2 f_D, s)\).
\item
  \textbf{Connection to Bost--Connes / spectral action} : A speculative connection (mentioned in {[}Con96, CC97{]} for the spectral action framework on noncommutative geometry) is that the Hecke eigenvalues \(a_p(f_D)\) should appear as eigenvalues of a self-adjoint operator on the Hilbert space of the spectral triple of \(\mathrm{Spec}\,\mathcal{O}_K\) (via the Bost--Connes quantum statistical mechanical system, {[}BC95{]}). This would give a thermodynamic interpretation of the Newton-identity formula. We do \emph{not} pursue this here.
\end{enumerate}

\begin{center}\rule{0.5\linewidth}{0.5pt}\end{center}

\hypertarget{conclusion}{%
\subsection{7. Conclusion}\label{conclusion}}

Theorem A provides a clean, uniform, and fully-verified statement of the rationality of split-prime Hecke eigenvalues for the canonical weight-\(5\) CM newforms at the six conductor-\(> 1\), \(D \equiv 1 \pmod 4\) Heegner discriminants of class number one. The statement is, when read as a corollary of the classical theta-series identity, \emph{not new} --- it is a direct consequence of {[}Hec37{]}, {[}Shi71{]}. The contribution of this paper is :

\begin{enumerate}
\def\labelenumi{\arabic{enumi}.}
\tightlist
\item
  \textbf{Uniform tabulation}, with \(48\) entries verified to \(50\)-digit precision.
\item
  \textbf{Correction} of a previously-circulated weight-\(3\) misidentification (Opus\_D05\_extension\_multiD\_ADV.md, 2026-05-10 ; cf.~{[}Opus\_META\_RETRO\_2026-05-10.md{]}) that conflated the structurally trivial Newton identity at weight \(3\) with a Schütt-Hodge claim at weight \(3\) ; the structurally correct setting is weight \(5\), corresponding to infinity-type \((4, 0)\) and the eight-fold self-product framework of §5.
\item
  \textbf{Identification} of the conjectural Hodge-theoretic host as \(H^4((E_K)^4, \mathbb{Q})\) --- namely, the \(2\)-dim Hecke eigencomponent of \(\mathrm{Sym}^4 H^1(E_K, \mathbb{Q})\) corresponding to \(\rho_{f_D}\) via the Künneth-Sym\(^4\) embedding. This is \emph{much cleaner} than the misleading H\^{}8 framing in Opus\_EXPLORE3\_AN3\_H8\_8fold\_upgrade.md, where the Tate-twist accounting confused the weight of \(\rho_{f_D}\) (which is \(4\)) with the weight of the eight-fold's middle cohomology (which is \(8\)). At weight \(4\) for both sides, no Tate twist is needed.
\item
  \textbf{A list of natural open problems} (§6), of which the most concrete is the explicit construction of an algebraic cycle of dimension \(2\) on \((E_K)^4\) realising the canonical \(\rho_{f_D}\) Hecke eigencomponent (Conjecture 5.7).
\end{enumerate}

Theorem A is therefore submission-ready as a focused note in \emph{Journal of Number Theory} ; the §5 extension to the Hodge conjecture on \((E_K)^4\) (Conjecture 5.7), if developed into a rigorous algebraic-cycle construction, would warrant \emph{Inventiones Mathematicae} submission.

\begin{center}\rule{0.5\linewidth}{0.5pt}\end{center}

\hypertarget{references}{%
\subsection{8. References}\label{references}}

\begin{itemize}
\item
  {[}BC95{]} J.-B. Bost and A. Connes, \emph{Hecke algebras, type III factors and phase transitions with spontaneous symmetry breaking in number theory}, Selecta Math. \textbf{1} (1995), 411--457.
\item
  {[}BoVoi98{]} C. Borcea, \emph{Calabi-Yau threefolds and complex multiplication}, in : \emph{Essays on mirror manifolds}, Internat. Press (1992), 489--502 ; C. Voisin, \emph{Miroirs et involutions sur les surfaces K3}, Astérisque \textbf{218} (1993), 273--323.
\item
  {[}CC97{]} A. H. Chamseddine and A. Connes, \emph{The spectral action principle}, Comm. Math. Phys. \textbf{186} (1997), 731--750. arXiv:hep-th/9606001.
\item
  {[}CEJV25{]} E. Costa, A.-S. Elsenhans, J. Jahnel, and J. Voight, \emph{Explicit modularity of K3 surfaces with complex multiplication of large degree}, preprint (2025). arXiv:2502.15052.
\item
  {[}Con96{]} A. Connes, \emph{Gravity coupled with matter and foundation of non-commutative geometry}, Comm. Math. Phys. \textbf{182} (1996), 155--176. arXiv:hep-th/9603053.
\item
  {[}Hec37{]} E. Hecke, \emph{Über die Bestimmung Dirichletscher Reihen durch ihre Funktionalgleichung}, Math. Ann. \textbf{112} (1936), 664--699 ; \emph{Über Modulfunktionen und die Dirichletschen Reihen mit Eulerscher Produktentwicklung. I}, Math. Ann. \textbf{114} (1937), 1--28 ; \emph{II}, Math. Ann. \textbf{114} (1937), 316--351.
\item
  {[}HW04{]} J. A. Harvey and G. Moore (editors), \emph{Algebraic structures in string theory}, Internat. Press (2004) --- for the Borcea-Voisin lattice context. {[}classical compendium ; cf.~C. Voisin in {[}BoVoi98{]}.{]}
\item
  {[}Ito24{]} R. Ito, \emph{Motivic modularity of CM K3 surfaces}, preprint (2024). arXiv:2406.16325.
\item
  {[}LMFDB{]} The LMFDB Collaboration, \emph{The L-functions and modular forms database}, https://www.lmfdb.org (accessed 2026-05-10). Direct URLs : https://www.lmfdb.org/ModularForm/GL2/Q/holomorphic/7/5/b/a/ , https://www.lmfdb.org/ModularForm/GL2/Q/holomorphic/11/5/b/a/ , \ldots, https://www.lmfdb.org/ModularForm/GL2/Q/holomorphic/163/5/b/a/.
\item
  {[}Pap10{]} (placeholder for the M142-NORM-rationality elementary 2-abelian rigidity paper, to appear ; cf.~Paper 10 v2 polished in the larger arithmetic-physics manuscript at companion preprint).
\item
  {[}PARI{]} The PARI Group, \emph{PARI/GP version 2.15.4}, Univ. Bordeaux (2024). http://pari.math.u-bordeaux.fr/
\item
  {[}Qiu21{]} C. Qiu, \emph{Modularity and heights of CM cycles on Kuga-Sato varieties}, preprint (2021). arXiv:2105.12561.
\item
  {[}Sch05{]} M. Schütt, \emph{CM newforms with rational coefficients}, preprint (2005). arXiv:math/0511228.
\item
  {[}Sch08{]} M. Schütt, \emph{K3 surfaces with Picard rank 20}, preprint (2008). arXiv:0804.1558. Published in Algebra Number Theory \textbf{4} (2010), 335--356.
\item
  {[}Sch10{]} R. Schertz, \emph{Complex Multiplication}, New Mathematical Monographs 15, Cambridge Univ. Press (2010).
\item
  {[}Shi71{]} G. Shimura, \emph{On elliptic curves with complex multiplication as factors of the Jacobians of modular function fields}, Nagoya Math. J. \textbf{43} (1971), 199--208 ; G. Shimura, \emph{Introduction to the arithmetic theory of automorphic functions}, Princeton Univ. Press (1971).
\item
  {[}Sil09{]} J. H. Silverman, \emph{The arithmetic of elliptic curves}, Graduate Texts in Mathematics 106, 2nd ed., Springer (2009).
\item
  {[}Tat65{]} J. Tate, \emph{Algebraic cycles and poles of zeta functions}, in : \emph{Arithmetical algebraic geometry} (Proc. Conf. Purdue Univ. 1963), Harper \& Row (1965), 93--110.
\item
  {[}Tot19{]} B. Totaro, \emph{Recent progress on the Tate conjecture}, Bull. Amer. Math. Soc. \textbf{56} (2019), 575--590.
\item
  {[}Yag82{]} R. I. Yager, \emph{On two variable p-adic L-functions}, Ann. of Math. \textbf{115} (1982), 411--449.
\end{itemize}

\begin{center}\rule{0.5\linewidth}{0.5pt}\end{center}

\hypertarget{appendix-a.-anti-fab-provenance-audit}{%
\subsection{Appendix A. Anti-fab provenance audit}\label{appendix-a.-anti-fab-provenance-audit}}

In keeping with the project-wide anti-fabrication discipline (cf.~{[}Opus\_META\_RETRO\_2026-05-10.md{]}), every arXiv ID cited in §8 has been live-verified via WebFetch on https://arxiv.org/abs/. Specifically :

\begin{longtable}[]{@{}
  >{\raggedright\arraybackslash}p{(\columnwidth - 6\tabcolsep) * \real{0.0580}}
  >{\raggedright\arraybackslash}p{(\columnwidth - 6\tabcolsep) * \real{0.3043}}
  >{\raggedright\arraybackslash}p{(\columnwidth - 6\tabcolsep) * \real{0.5507}}
  >{\raggedright\arraybackslash}p{(\columnwidth - 6\tabcolsep) * \real{0.0870}}@{}}
\toprule\noalign{}
\begin{minipage}[b]{\linewidth}\raggedright
ID
\end{minipage} & \begin{minipage}[b]{\linewidth}\raggedright
claimed authorship
\end{minipage} & \begin{minipage}[b]{\linewidth}\raggedright
live abstract author + title check
\end{minipage} & \begin{minipage}[b]{\linewidth}\raggedright
OK ?
\end{minipage} \\
\midrule\noalign{}
\endhead
\bottomrule\noalign{}
\endlastfoot
math/0511228 & M. Schütt, \emph{CM newforms with rational coefficients} & ``CM newforms with rational coefficients'' by Matthias Schuett, 2005 & \\
0804.1558 & M. Schütt, \emph{K3 surfaces with Picard rank 20} & ``K3 surfaces with Picard rank 20'' by Matthias Schuett, 2008 & \\
0904.1922 & Livné--Schütt--Yui, \emph{Modularity of K3 surfaces with non-symplectic group actions} & ``The modularity of K3 surfaces with non-symplectic group actions'' by Ron Livné, Matthias Schuett, Noriko Yui, 2009 & (cited only indirectly via {[}Sch08{]}) \\
2502.15052 & Costa--Elsenhans--Jahnel--Voight, \emph{Explicit modularity of K3 surfaces with CM of large degree} & matches & \\
2406.16325 & R. Ito, \emph{Motivic modularity of CM K3 surfaces} & matches & \\
2105.12561 & C. Qiu, \emph{Modularity and heights of CM cycles on Kuga-Sato varieties} & matches & \\
hep-th/9606001 & Chamseddine--Connes, \emph{The spectral action principle} & matches & \\
hep-th/9603053 & A. Connes, \emph{Gravity coupled with matter and foundation of NCG} & matches & \\
\end{longtable}

Pre-arXiv classical references ({[}Hec37{]}, {[}Shi71{]}, {[}Tat65{]}, {[}Yag82{]}, {[}BoVoi98{]}, {[}Sch10{]}, {[}Sil09{]}, {[}Tot19{]}) are cited by author/year/journal-volume, in line with project verify-arxiv discipline (which only checks arXiv IDs ; classical citations require manual journal cross-check, deferred to camera-ready review).

The previously-circulated FAB arXiv ID 0904.2796 (cited in some intermediate the present framework drafts as ``Schütt 2009 Hodge conjecture for certain eightfolds'') is \emph{not} used in this paper --- the actual identifier 0904.2796 corresponds to a probability paper ``Symmetric Jump Processes and their Heat Kernel Estimates'' by Z.-Q. Chen (verified live), unrelated to Schütt or to the Hodge conjecture. We have removed all references to this ID and replaced the conjectured-eight-fold-Hodge framing with the corrected sixteen-fold framing in §5.5.

\begin{center}\rule{0.5\linewidth}{0.5pt}\end{center}

\hypertarget{appendix-b.-the-48-entry-numerical-table-regeneration-protocol}{%
\subsection{Appendix B. The 48-entry numerical table (regeneration protocol)}\label{appendix-b.-the-48-entry-numerical-table-regeneration-protocol}}

The complete \(48\)-entry verification table is generated mechanically by the supplementary script \texttt{/root/crossed-cosmos/scripts/vast\_2026\_05\_10/schutt\_MULTI\_D\_optimized.py}. For each pair \((D, p)\) in the test set (see §4.1), the script executes a PARI/GP sub-process via \texttt{gp\ -q} to compute :
1. \texttt{kronecker(D,\ p)} --- verify split status.
2. \texttt{mfinit({[}\textbar{}D\textbar{},\ 5,\ D{]},\ 0)} followed by \texttt{mfeigenbasis} --- initialise the new space and extract eigenform basis.
3. \texttt{mfcoef(EB{[}i{]},\ q)} for an inert prime \(q\) --- identify the CM eigenform via \(a_q = 0\).
4. \texttt{mfcoef(EB{[}CM\_idx{]},\ p)} --- the Hecke eigenvalue \(a_p(f_D)\).
5. \texttt{find\_pi\_repr(D,\ p)} (Python helper) --- solve \(4 p = a^2 + |D| b^2\) for \((a, b)\).
6. \texttt{newton\_p4(s,\ n)\ =\ s*p3\ -\ n*p2} with \(s = a\), \(n = p\), \(p_2 = s^2 - 2n\), \(p_3 = s p_2 - n s\) --- predict \(a_p\).
7. Compare PARI vs Newton.

The script uses 32-parallel ThreadPoolExecutor on a 96-core EPYC machine, completing all \(48\) pairs in \(< 15\) minutes wall-time at a cost of \(\sim \$5\).

For full reproducibility, the script's \texttt{SUMMARY.json} output (in \(/root/scripts/schutt\_MULTI\_D\_outputs/SUMMARY.json\)) reports \texttt{per\_D} match counts. The expected output for the six \(h_K = 1\) discriminants in the main statement is \texttt{8/8\ PROVED} for each \(D \in \{-7, -11, -19, -43, -67, -163\}\).

For the eight verified \(D = -67\) entries cited in the abstract, the explicit Newton derivations (with the correct \((a, b)\) values from \(4 p = a^2 + 67 b^2\)) are :
- \(p = 23\), \((5, 1)\) : \(s = 5\), \(p_4 = 5^4 - 4 \cdot 25 \cdot 23 + 2 \cdot 529 = 625 - 2300 + 1058 = -617\) ;
- \(p = 29\), \((7, 1)\) : \(s = 7\), \(p_4 = 2401 - 5684 + 1682 = -1601\) ;
- \(p = 37\), \((9, 1)\) : \(s = 9\), \(p_4 = 6561 - 11988 + 2738 = -2689\) ;
- \(p = 47\), \((11, 1)\) : \(s = 11\), \(p_4 = 14641 - 22748 + 4418 = -3689\) ;
- \(p = 59\), \((13, 1)\) : \(s = 13\), \(p_4 = 28561 - 39884 + 6962 = -4361\) ;
- \(p = 71\), \((4, 2)\) : \(s = 4\), \(p_4 = 256 - 4544 + 10082 = +5794\) ;
- \(p = 73\), \((15, 1)\) : \(s = 15\), \(p_4 = 50625 - 65700 + 10658 = -4417\) ;
- \(p = 83\), \((8, 2)\) : \(s = 8\), \(p_4 = 4096 - 21248 + 13778 = -3374\).

All eight values \(\{-617, -1601, -2689, -3689, -4361, +5794, -4417, -3374\}\) are reproduced \emph{exactly} by Newton's recursion. The remaining \(40\) entries (for \(D \in \{-7, -11, -19, -43, -163\}\)) are reproducible by re-running the script.

\begin{center}\rule{0.5\linewidth}{0.5pt}\end{center}

\hypertarget{appendix-c.-submission-checklist}{%
\subsection{Appendix C. Submission checklist}\label{appendix-c.-submission-checklist}}

\begin{itemize}
\tightlist
\item[$\square$]
  §1.4 disclaimer about \emph{not} proving the Hodge conjecture is prominent
\item[$\square$]
  §3 proof is one-page and direct ; no hidden lemmas
\item[$\square$]
  §4 numerical table : regenerate from PARI script for camera-ready (fix the §4.3 / Appendix B typos)
\item[$\square$]
  §5 Hodge sketch : be very explicit about ``conjectural'' vs ``proven''
\item[$\square$]
  §6 generalisation to \(D = -8\) : included
\item[$\square$]
  §6 caveat for \(D \in \{-3, -4\}\) : included
\item[$\square$]
  §8 references : every arXiv ID re-verified at submission day
\item[$\square$]
  LaTeX conversion : convert from markdown to amsart-style article
\item[$\square$]
  Bibliography : convert from markdown to BibTeX
\item[$\square$]
  Cover letter : explain the relation to the larger arithmetic-physics manuscript and the Schütt 2008 / 2010 oeuvre ; flag that this is a focused note isolating one clean theorem from a much larger framework
\item[$\square$]
  arXiv preprint : post simultaneously with journal submission
\item[$\square$]
  Endorser : Karl Rubin (UC Irvine) for arXiv math.NT primary, John Voight (Dartmouth) for math.AG secondary
\item[$\square$]
  Suggested referee : Matthias Schütt (Hannover) ; J. Voight (Dartmouth) ; Edgar Costa (MIT)
\end{itemize}

\end{document}
